% Pertemuan 5 Praktikum Pemrograman (IF25-11002). Dihasilkan oleh tools/build.py dari tools/konten/p05.py.
% Kompilasi: xelatex P05.tex (dua kali). Catatan: xelatex -jobname=P05-catatan "\def\catatan{1}% Pertemuan 5 Praktikum Pemrograman (IF25-11002). Dihasilkan oleh tools/build.py dari tools/konten/p05.py.
% Kompilasi: xelatex P05.tex (dua kali). Catatan: xelatex -jobname=P05-catatan "\def\catatan{1}% Pertemuan 5 Praktikum Pemrograman (IF25-11002). Dihasilkan oleh tools/build.py dari tools/konten/p05.py.
% Kompilasi: xelatex P05.tex (dua kali). Catatan: xelatex -jobname=P05-catatan "\def\catatan{1}% Pertemuan 5 Praktikum Pemrograman (IF25-11002). Dihasilkan oleh tools/build.py dari tools/konten/p05.py.
% Kompilasi: xelatex P05.tex (dua kali). Catatan: xelatex -jobname=P05-catatan "\def\catatan{1}\input{P05.tex}"
\documentclass[t]{beamer}
\usetheme{ITERA}
\input{catatan-preamble.tex}
\renewcommand\ITERAmeeting{Pertemuan 5}
\begin{document}
\SlideJudul{IF25-11002 PRAKTIKUM PEMROGRAMAN}{Pertemuan 5\\Perulangan Bagian 1}{for, while, dan do-while dengan pola pencacah, akumulator, dan sentinel}{Tim Pengajar}{Tahun 2026. Sejajar dengan Pertemuan 5 Algoritma Pemrograman: Perulangan Bagian 1.}
\note{Buka dengan menanyakan satu hal dari pertemuan lalu yang masih membingungkan.}

\begin{frame}{Dua mata kuliah, satu minggu}
\framesubtitle{Sebelum mulai}
Topik minggu ini sama di dua mata kuliah, dengan peran yang berbeda.
\vspace{10pt}

\begin{columns}[T]
\begin{column}{0.47\textwidth}
\begin{Blok}{Algoritma: Perulangan Bagian 1}
Memodelkan pengulangan dengan flowchart dan pseudocode ULANGI dan SELAMA, lalu menelusurinya dengan trace table.
\end{Blok}
\end{column}
\begin{column}{0.47\textwidth}
\begin{BlokGelap}{Praktikum: Perulangan 1}
Menulis \texttt{for}, \texttt{while}, dan \texttt{do-while}, memakai pencacah dan akumulator, dan menghentikan perulangan dengan sentinel.
\end{BlokGelap}
\end{column}
\end{columns}
\note{Tanyakan apa yang sudah dibahas di Algoritma minggu ini. Gunakan jawabannya sebagai jembatan ke praktikum.}
\end{frame}

\begin{frame}{Saat pulang nanti, kalian bisa}
\framesubtitle{Target hari ini}
\begin{columns}[T]
\begin{column}{0.47\textwidth}
\begin{Langkah}{1}{Menerapkan}
Menulis program yang mengulang proses dengan \texttt{for}, \texttt{while}, dan \texttt{do-while}, memakai pola pencacah, akumulator, dan nilai sentinel

{\footnotesize\color{ITERAInkMuted}Sub-CPMK 5.1, CPMK0607}
\end{Langkah}
\end{column}
\begin{column}{0.47\textwidth}
\begin{Langkah}{2}{Menjelaskan}
Menelusuri perulangan dengan trace table, menentukan berapa kali badan perulangan dijalankan, dan menjelaskan penyebab off-by-one atau perulangan tak berhenti

{\footnotesize\color{ITERAInkMuted}Sub-CPMK 5.2, CPMK0608}
\end{Langkah}
\end{column}
\end{columns}
\vspace{8pt}
\begin{Catatan}
Dua kemampuan ini dinilai sama berat di Exit Ticket, tugas, UTS, dan UAS.
\end{Catatan}
\note{Bacakan target dengan pelan. Kata kuncinya menerapkan dan menjelaskan.}
\end{frame}

\begin{frame}{Ingat minggu lalu}
\framesubtitle{Review, 5 menit}
\begin{columns}[T]
\begin{column}{0.31\textwidth}
\begin{BlokGelap}{Pertanyaan 1}
Apa keluaran rantai \texttt{else if} untuk nilai batas 75?
\end{BlokGelap}
\end{column}
\begin{column}{0.31\textwidth}
\begin{BlokGelap}{Pertanyaan 2}
Mengapa perlu \texttt{break} di \texttt{switch}?
\end{BlokGelap}
\end{column}
\begin{column}{0.31\textwidth}
\begin{BlokGelap}{Pertanyaan 3}
Apa beda \texttt{if (x = 5)} dan \texttt{if (x == 5)}?
\end{BlokGelap}
\end{column}
\end{columns}
\vspace{8pt}
Jawab di kertas dulu, lalu cocokkan dengan teman sebelah.\note{Pilih dua mahasiswa untuk menjawab. Koreksi singkat, jangan mengajar ulang.}
\end{frame}

\Bagian{1}{Masalah pembuka}{Rata-rata nilai satu kelas}
\note{Masalah pembuka 10 menit.}

\begin{frame}[fragile]{Rata-rata nilai satu kelas}
\framesubtitle{Masalah minggu ini}
\begin{columns}[T]
\begin{column}{0.55\textwidth}
{Asisten praktikum harus menghitung rata-rata nilai Exit Ticket. Jumlah mahasiswa berbeda di setiap kelas, kadang 30, kadang 45.

\vspace{6pt}
Menulis \texttt{cin} sebanyak 45 kali jelas tidak masuk akal. Lagi pula banyaknya data baru diketahui saat program berjalan.\par}
\vspace{6pt}
\begin{Catatan}
\textbf{Diskusi 2 menit.} Langkah apa yang terus diulang untuk setiap nilai? Bagaimana program tahu kapan berhenti?
\end{Catatan}
\end{column}
\begin{column}{0.41\textwidth}
\begin{Keluaran}[]
Masukkan nilai (-1 untuk selesai):
70 85 60 90 -1
Banyak    : 4
Tertinggi : 90
Rata-rata : 76.25
\end{Keluaran}
\end{column}
\end{columns}
\note{Tampung beberapa jawaban tanpa menilai benar salah.}
\end{frame}

\begin{frame}{Dari langkah ke instruksi}
\framesubtitle{Sambungan dengan Algoritma}
\begin{columns}[T]
\begin{column}{0.31\textwidth}
\begin{Langkah}{1}{Temukan yang berulang}
Baca satu nilai, tambahkan ke total, tambah pencacah.
\end{Langkah}
\end{column}
\begin{column}{0.31\textwidth}
\begin{Langkah}{2}{Tentukan kapan berhenti}
Berhenti saat bertemu -1, atau setelah n kali.
\end{Langkah}
\end{column}
\begin{column}{0.31\textwidth}
\begin{Langkah}{3}{Terjemahkan}
\texttt{while} untuk syarat berhenti, \texttt{for} untuk banyak putaran yang sudah diketahui.
\end{Langkah}
\end{column}
\end{columns}
\vspace{10pt}
Langkah 1 dan 2 dilatih di Algoritma. Langkah 3 kita kerjakan hari ini dengan C++.\note{Hubungkan eksplisit ke Algoritma minggu ini.}
\end{frame}

\Bagian{2}{Coba dulu, baru dijelaskan}{25 menit eksperimen di GDB Online}
\note{Struggle 25 menit. Berkeliling, jangan menjelaskan materi dulu.}

\begin{frame}{Misi kalian}
\framesubtitle{Struggle, 25 menit}
\begin{columns}[T]
\begin{column}{0.31\textwidth}
\begin{Langkah}{1}{Hitung}
Tampilkan angka 1 sampai 10 tanpa menulis \texttt{cout} sepuluh kali.
\end{Langkah}
\end{column}
\begin{column}{0.31\textwidth}
\begin{Langkah}{2}{Jumlahkan}
Hitung 1 + 2 + ... + 100. Tebak dulu hasilnya.
\end{Langkah}
\end{column}
\begin{column}{0.31\textwidth}
\begin{Langkah}{3}{Berhenti}
Baca angka terus sampai pengguna mengetik 0.
\end{Langkah}
\end{column}
\end{columns}
\vspace{8pt}
\begin{columns}[T]
\begin{column}{0.47\textwidth}
\begin{Blok}{Boleh}
Bertanya ke teman, mengubah apa saja, sengaja membuat error.
\end{Blok}
\end{column}
\begin{column}{0.47\textwidth}
\begin{BlokAlert}{Belum dulu}
Mencari jawaban jadi di internet atau AI.
\end{BlokAlert}
\end{column}
\end{columns}
\note{Catat kesalahan yang sering muncul untuk dibahas di debrief.}
\end{frame}

\begin{frame}{Apa yang kalian temukan?}
\framesubtitle{Debrief struggle}
\begin{columns}[T]
\begin{column}{0.31\textwidth}
\begin{BlokGelap}{Pertanyaan 1}
Variabel apa yang kalian butuhkan untuk menyimpan jumlah?
\end{BlokGelap}
\end{column}
\begin{column}{0.31\textwidth}
\begin{BlokGelap}{Pertanyaan 2}
Di mana variabel itu harus diberi nilai awal?
\end{BlokGelap}
\end{column}
\begin{column}{0.31\textwidth}
\begin{BlokGelap}{Pertanyaan 3}
Apakah ada yang programnya tidak mau berhenti? Mengapa?
\end{BlokGelap}
\end{column}
\end{columns}
\vspace{10pt}
Jawaban kalian akan kita uji di bagian berikutnya.\note{Tulis jawaban di papan; buktikan atau patahkan di bagian materi.}
\end{frame}

\Bagian{3}{Mengulang dengan terkendali}{for, while, do-while, dan pola-pola perulangan}
\note{Materi 40 menit. Tunjukkan setiap contoh langsung di GDB Online.}

\begin{frame}[fragile]{Perulangan for}
\framesubtitle{Banyak putaran diketahui}
\begin{columns}[T]
\begin{column}{0.55\textwidth}
\begin{Kode}[]{for.cpp}
for (int i = 1; i <= 5; i++) {
    cout << i << " ";
}
cout << endl;
for (int i = 10; i > 0; i -= 3) {
    cout << i << " ";
}
cout << endl;
\end{Kode}
\end{column}
\begin{column}{0.41\textwidth}
\only<1>{\begin{TebakDulu}
{\ITERAKickerFont\fontsize{10pt}{12pt}\selectfont\color{ITERAAlert}TEBAK DULU}\par\vspace{4pt}
Sebelum dijalankan, tulis keluarannya di kertas.
\end{TebakDulu}
}\only<2>{\KeluaranFile[]{keluaran/P05/k001.txt}
\begin{Catatan}
Tiga bagian \texttt{for}: nilai awal; syarat lanjut; perubahan. Syarat diperiksa sebelum setiap putaran.
\end{Catatan}
}
\end{column}
\end{columns}
\note<1>{Minta mahasiswa menebak keluaran sebelum membuka slide berikutnya. Tanyakan satu atau dua tebakan yang berbeda, lalu buktikan.}
\end{frame}

\begin{frame}[fragile]{Pencacah dan akumulator}
\framesubtitle{Pola perulangan}
\begin{columns}[T]
\begin{column}{0.55\textwidth}
\begin{Kode}[]{akumulator.cpp}
int jumlah = 0, genap = 0;
for (int i = 1; i <= 10; i++) {
    jumlah += i;
    if (i % 2 == 0) {
        genap++;
    }
}
cout << "Jumlah: " << jumlah << endl;
cout << "Genap : " << genap << endl;
\end{Kode}
\end{column}
\begin{column}{0.41\textwidth}
\only<1>{\begin{TebakDulu}
{\ITERAKickerFont\fontsize{10pt}{12pt}\selectfont\color{ITERAAlert}TEBAK DULU}\par\vspace{4pt}
Sebelum dijalankan, tulis keluarannya di kertas.
\end{TebakDulu}
}\only<2>{\KeluaranFile[]{keluaran/P05/k002.txt}
\begin{Catatan}
Akumulator dan pencacah selalu diberi nilai awal sebelum perulangan, bukan di dalamnya.
\end{Catatan}
}
\end{column}
\end{columns}
\note<1>{Tanya: apa yang terjadi bila int jumlah = 0 dipindah ke dalam perulangan?}
\end{frame}

\begin{frame}[fragile]{Perulangan while}
\framesubtitle{Berhenti karena syarat}
\begin{columns}[T]
\begin{column}{0.60\textwidth}
\begin{Kode}[listing options={style=ITERACodeKecil}]{while.cpp}
int saldo = 100000, bulan = 0;
while (saldo < 500000) {
    saldo += 150000;
    bulan++;
}
cout << bulan << " bulan, saldo " << saldo << endl;
\end{Kode}
\end{column}
\begin{column}{0.36\textwidth}
\only<1>{\begin{TebakDulu}
{\ITERAKickerFont\fontsize{10pt}{12pt}\selectfont\color{ITERAAlert}TEBAK DULU}\par\vspace{4pt}
Sebelum dijalankan, tulis keluarannya di kertas.
\end{TebakDulu}
}\only<2>{\KeluaranFile[listing options={style=ITERAPlainKecil}]{keluaran/P05/k003.txt}
\begin{Catatan}
Sesuatu di dalam badan harus mengubah syarat. Bila tidak, perulangan tidak pernah berhenti.
\end{Catatan}
}
\end{column}
\end{columns}
\note<1>{Minta mahasiswa menebak keluaran sebelum membuka slide berikutnya. Tanyakan satu atau dua tebakan yang berbeda, lalu buktikan.}
\end{frame}

\begin{frame}[fragile]{Nilai sentinel}
\framesubtitle{Membaca sampai tanda berhenti}
\begin{columns}[T]
\begin{column}{0.60\textwidth}
\begin{Kode}[listing options={style=ITERACodeKecil}]{sentinel.cpp}
int nilai, banyak = 0, total = 0;
cin >> nilai;
while (nilai != -1) {
    banyak++;
    total += nilai;
    cin >> nilai;
}
cout << banyak << " data, total " << total << endl;
\end{Kode}
\end{column}
\begin{column}{0.36\textwidth}
\begin{Keluaran}[title={Masukan},listing options={style=ITERAPlainKecil}]
70 85 60 90 -1
\end{Keluaran}
\only<1>{\begin{TebakDulu}
{\ITERAKickerFont\fontsize{10pt}{12pt}\selectfont\color{ITERAAlert}TEBAK DULU}\par\vspace{4pt}
Sebelum dijalankan, tulis keluarannya di kertas.
\end{TebakDulu}
}\only<2>{\KeluaranFile[listing options={style=ITERAPlainKecil}]{keluaran/P05/k004.txt}
\begin{Catatan}
Baca sekali sebelum perulangan, lalu baca lagi di akhir badan. Dengan begitu -1 tidak ikut dihitung.
\end{Catatan}
}
\end{column}
\end{columns}
\note<1>{Minta mahasiswa menebak keluaran sebelum membuka slide berikutnya. Tanyakan satu atau dua tebakan yang berbeda, lalu buktikan.}
\end{frame}

\begin{frame}[fragile]{Perulangan do-while}
\framesubtitle{Minimal sekali jalan}
\begin{columns}[T]
\begin{column}{0.60\textwidth}
\begin{Kode}[listing options={style=ITERACodeKecil}]{dowhile.cpp}
int pin, coba = 0;
do {
    cin >> pin;
    coba++;
} while (pin != 1234 && coba < 3);
cout << (pin == 1234 ? "Masuk" : "Terkunci");
cout << " setelah " << coba << " kali" << endl;
\end{Kode}
\end{column}
\begin{column}{0.36\textwidth}
\begin{Keluaran}[title={Masukan},listing options={style=ITERAPlainKecil}]
1111 1234
\end{Keluaran}
\only<1>{\begin{TebakDulu}
{\ITERAKickerFont\fontsize{10pt}{12pt}\selectfont\color{ITERAAlert}TEBAK DULU}\par\vspace{4pt}
Sebelum dijalankan, tulis keluarannya di kertas.
\end{TebakDulu}
}\only<2>{\KeluaranFile[listing options={style=ITERAPlainKecil}]{keluaran/P05/k005.txt}
\begin{Catatan}
Syarat diperiksa sesudah badan. Jangan lupa titik koma setelah \texttt{while (...)}.
\end{Catatan}
}
\end{column}
\end{columns}
\note<1>{Minta mahasiswa menebak keluaran sebelum membuka slide berikutnya. Tanyakan satu atau dua tebakan yang berbeda, lalu buktikan.}
\end{frame}

\begin{frame}{Memilih jenis perulangan}
\framesubtitle{for, while, atau do-while}
\small
\begin{tabular}{@{}>{\raggedright\arraybackslash}p{364pt}>{\raggedright\arraybackslash}p{153pt}>{\raggedright\arraybackslash}p{345pt}@{}}
\kepala{Situasi} & \kepala{Pilihan} & \kepala{Contoh}\\
Banyak putaran diketahui & \texttt{for} & tampilkan 1 sampai n\\
\barisgenap Berhenti karena syarat, bisa nol putaran & \texttt{while} & baca sampai sentinel\\
Badan harus jalan minimal sekali & \texttt{do-while} & menu, validasi PIN\\
\garisdasar
\end{tabular}
\vspace{10pt}

Ketiganya bisa saling menggantikan. Pilih yang paling jelas menyatakan maksud kalian.
\end{frame}

\begin{frame}[fragile]{Tebak keluarannya}
\framesubtitle{Latihan menjelaskan, 3 menit}
\begin{columns}[T]
\begin{column}{0.56\textwidth}
\begin{Kode}[]{tebak.cpp}
int x = 0;
for (int i = 0; i < 4; i++) {
    x += i * 2;
}
int k = 1;
while (k < 20) {
    k *= 3;
}
cout << x << " " << k << endl;
\end{Kode}
\end{column}
\begin{column}{0.40\textwidth}
\begin{Catatan}
Tulis keluarannya di kertas, karakter demi karakter. Jangan dijalankan dulu.
\end{Catatan}
\end{column}
\end{columns}
\note{Contoh soal Bagian B. Beri waktu 3 menit sebelum slide jawaban.}
\end{frame}

\begin{frame}[fragile]{Tebak keluarannya: jawaban}
\framesubtitle{Latihan menjelaskan}
\begin{columns}[T]
\begin{column}{0.56\textwidth}
\begin{Kode}[]{tebak.cpp}
int x = 0;
for (int i = 0; i < 4; i++) {
    x += i * 2;
}
int k = 1;
while (k < 20) {
    k *= 3;
}
cout << x << " " << k << endl;
\end{Kode}
\end{column}
\begin{column}{0.40\textwidth}
\begin{Keluaran}[]
12 27
\end{Keluaran}
\begin{Catatan}
\texttt{x} = 0 + 2 + 4 + 6 = 12. \texttt{k} menjadi 3, 9, 27; berhenti saat 27 tidak lagi kurang dari 20.
\end{Catatan}
\end{column}
\end{columns}
\note{Tanyakan jawaban yang paling sering salah, lalu telusuri bersama.}
\end{frame}

\begin{frame}{Error yang sering muncul minggu ini}
\framesubtitle{Pesan diambil langsung dari g++}
\footnotesize
\begin{tabular}{@{}>{\raggedright\arraybackslash}p{208pt}>{\raggedright\arraybackslash}p{256pt}>{\raggedright\arraybackslash}p{398pt}@{}}
\kepala{Kesalahan} & \kepala{Contoh} & \kepala{Pesan pertama dari g++}\\
Koma di header for & \texttt{for (int i = 0, i < 5, i++)} & \texttt{error: expected ';' before '<' token}\\
\barisgenap Variabel di luar jangkauan & \texttt{cout << i; sesudah for} & \texttt{error: 'i' was not declared in this scope}\\
Lupa ; sesudah do-while & \texttt{\} while (x < 5)} & \texttt{error: expected ';' before 'cout'}\\
\barisgenap Perbandingan bertipe beda & \texttt{i < s.length() dengan int} & \texttt{warning: comparison of integer expressions of different signedness: 'int' and ...}\\
Badan for kosong & \texttt{for (...);} & \texttt{warning: this 'for' clause does not guard...}\\
\garisdasar
\end{tabular}
\vspace{10pt}

{\small Pesan di tabel ini dihasilkan dengan g++ dan opsi -Wall, sama seperti di cek.sh.}
\note{Minta mahasiswa memotret slide ini.}
\end{frame}

\Bagian{4}{Hands-on bertingkat}{45 menit, dari dasar sampai tantangan}
\note{Mahasiswa membuka folder 02 Hands-on/P5 - Perulangan 1 - Hands-on.}

\begin{frame}{Peta hands-on}
\framesubtitle{Folder 02 Hands-on/P5 - Perulangan 1 - Hands-on}
\small
\begin{tabular}{@{}>{\raggedright\arraybackslash}p{36pt}>{\raggedright\arraybackslash}p{267pt}>{\raggedright\arraybackslash}p{110pt}>{\raggedright\arraybackslash}p{83pt}>{\raggedright\arraybackslash}p{341pt}@{}}
\kepala{No} & \kepala{File} & \kepala{Tingkat} & \kepala{Waktu} & \kepala{Yang dilatih}\\
1 & \texttt{handson1-deret.cpp} & Dasar & 10' & \texttt{for}, akumulator\\
\barisgenap 2 & \texttt{handson2-tabungan.cpp} & Menengah & 15' & \texttt{while} dengan syarat berhenti\\
3 & \texttt{handson3-rerata.cpp} & Tantangan & 20' & trace table, perbaiki tiga kesalahan sentinel\\
\barisgenap + & \texttt{bonus-faktorial.cpp} & Pengayaan & sisa & \texttt{do-while} untuk menu, \texttt{long long}\\
\garisdasar
\end{tabular}
\vspace{10pt}

Kerjakan berurutan. Cocokkan keluaran dengan folder \texttt{expected/}; latihan yang membaca masukan memakai data di folder \texttt{input/}.\note{Saat berkeliling, minta mahasiswa yang mengerjakan latihan tantangan menjelaskan blok PENJELASAN mereka.}
\end{frame}

\begin{frame}{Cara mengecek dan aturannya}
\framesubtitle{Hands-on}
\begin{columns}[T]
\begin{column}{0.47\textwidth}
\begin{Blok}{Di GDB Online}
Salin isi file latihan, lengkapi TODO, klik Run. Bila program membaca masukan, ketik isi file di \texttt{input/}. Bandingkan dengan \texttt{expected/}.
\end{Blok}
\end{column}
\begin{column}{0.47\textwidth}
\begin{BlokGelap}{Di komputer sendiri}
Jalankan \texttt{bash cek.sh}. Skrip mengompilasi dengan \texttt{-Wall -Werror}, memberi masukan dari \texttt{input/}, lalu mencocokkan keluaran.
\end{BlokGelap}
\end{column}
\end{columns}
\vspace{6pt}
\begin{Catatan}
AI boleh untuk memahami pesan error, tidak untuk meminta solusi utuh. Exit Ticket dikerjakan tanpa bantuan apa pun.
\end{Catatan}
\note{Hands-on tidak dinilai langsung; yang dinilai Exit Ticket dan tugas.}
\end{frame}

\Bagian{5}{Exit Ticket dan tugas}{25 menit terakhir, lalu pekerjaan rumah}
\note{Mulai Exit Ticket tepat waktu.}

\begin{frame}{Exit Ticket 5}
\framesubtitle{Individual, 25 menit, tanpa AI dan internet}
\small
\begin{tabular}{@{}>{\raggedright\arraybackslash}p{60pt}>{\raggedright\arraybackslash}p{560pt}>{\raggedright\arraybackslash}p{80pt}>{\raggedright\arraybackslash}p{150pt}@{}}
\kepala{Soal} & \kepala{Isi} & \kepala{Poin} & \kepala{CPMK}\\
A1 & Tulis program yang menampilkan kelipatan k dari 1 sampai n beserta banyaknya & 25 & CPMK0607\\
\barisgenap A2 & Tulis program yang membaca bilangan sampai 0 lalu menampilkan jumlah bilangan positif & 25 & CPMK0607\\
B1 & Buat trace table sebuah perulangan dan tentukan keluarannya & 20 & CPMK0608\\
\barisgenap B2 & Jelaskan mengapa sebuah perulangan tidak berhenti dan cara memperbaikinya & 20 & CPMK0608\\
B3 & Tentukan berapa kali badan perulangan dijalankan dan jelaskan alasannya & 10 & CPMK0608\\
\garisdasar
\end{tabular}
\vspace{10pt}

{\small Soal A dinilai dari keluaran yang tepat sama. Soal B dinilai dari ketepatan dan alasan. Pelanggaran aturan membuat nilai Exit Ticket ini 0.}
\note{Soal lengkap dibagikan terpisah dan tidak disimpan di repo publik.}
\end{frame}

\begin{frame}{Tugas 5: Program Statistik Nilai}
\framesubtitle{Dikumpulkan sebelum Pertemuan 6}
\begin{columns}[T]
\begin{column}{0.47\textwidth}
\begin{Blok}{Bagian A, program, 50 poin}
Buat program statistik nilai dengan ketentuan berikut. Ketentuannya di slide berikut.
\end{Blok}
\end{column}
\begin{column}{0.47\textwidth}
\begin{BlokGelap}{Bagian B, penjelasan, 50 poin}
Komentar maksimal 150 kata dan trace table untuk n = 3 dengan nilai 80, 70, 90. Jelaskan alasan memilih \texttt{for}, \texttt{while}, atau \texttt{do-while} di setiap bagian, dan satu kesalahan off-by-one atau perulangan tak berhenti yang kalian temui.
\end{BlokGelap}
\end{column}
\end{columns}
\vspace{8pt}
{\small Data di tugas adalah data kalian sendiri. Rubrik lengkap ada di Modul Belajar 5.}
\note{Ingatkan bahwa penjelasan disalin dari AI mudah dikenali karena tidak menyebut pengalaman mereka sendiri.}
\end{frame}

\begin{frame}{Ketentuan Tugas 5}
\framesubtitle{Program Statistik Nilai}
\small
\begin{tabular}{@{}>{\raggedright\arraybackslash}p{87pt}>{\raggedright\arraybackslash}p{788pt}@{}}
\kepala{No} & \kepala{Ketentuan Bagian A}\\
1 & Baca jumlah siswa n, ulangi sampai n > 0\\
\barisgenap 2 & Baca nilai n siswa; setiap nilai harus 0 sampai 100, ulangi bila tidak valid\\
3 & Tampilkan nilai tertinggi, nilai terendah, dan rata-rata kelas\\
\barisgenap 4 & Tampilkan jumlah siswa LULUS (nilai 75 ke atas) dan TIDAK LULUS\\
5 & Menu \texttt{do-while} untuk mengulang atau keluar\\
\barisgenap Bonus & Tampilkan persentase kelulusan dengan dua desimal.\\
\garisdasar
\end{tabular}
\note{Bacakan ketentuan satu per satu dan tanyakan apakah ada yang belum jelas.}
\end{frame}

\begin{frame}{Rubrik Tugas 5}
\framesubtitle{Empat level, sama di semua instrumen}
\footnotesize
\begin{tabular}{@{}>{\raggedright\arraybackslash}p{166pt}>{\raggedright\arraybackslash}p{44pt}>{\raggedright\arraybackslash}p{166pt}>{\raggedright\arraybackslash}p{156pt}>{\raggedright\arraybackslash}p{146pt}>{\raggedright\arraybackslash}p{146pt}@{}}
\kepala{Kriteria} & \kepala{Poin} & \kepala{Sangat baik 85--100} & \kepala{Baik 70--84} & \kepala{Cukup 55--69} & \kepala{Kurang <55}\\
A1 Program berjalan & 20 & tanpa error dan peringatan & ada peringatan & perlu satu perbaikan & tidak terkompilasi\\
\barisgenap A2 Validasi dan perulangan & 15 & validasi n dan nilai, menu berulang & satu validasi kurang & dua kurang & tanpa validasi\\
A3 Ketepatan statistik & 15 & semua tepat termasuk n = 1 & satu salah & dua salah & sebagian besar salah\\
\barisgenap B1 Trace table dan alasan & 30 & trace lengkap dan alasan tepat & satu baris trace salah & trace tidak lengkap & tidak ada atau disalin\\
B2 Cerita error dan pengujian & 20 & pesan, sebab, perbaikan, uji & pesan tidak dikutip & hanya perbaikan & tidak ada\\
\garisdasar
\end{tabular}
\vspace{8pt}

{\small Nilai kriteria = poin × skor level. Bagian A total 50, Bagian B total 50.}
\note{Rubrik ini sama strukturnya setiap minggu; yang berubah hanya isi kriteria A2, A3, dan B1.}
\end{frame}

\begin{frame}{Sebelum Pertemuan 6}
\framesubtitle{Persiapan}
\begin{columns}[T]
\begin{column}{0.31\textwidth}
\begin{Langkah}{1}{Selesaikan}
Hands-on yang belum lulus \texttt{cek.sh}, terutama latihan tantangan.
\end{Langkah}
\end{column}
\begin{column}{0.31\textwidth}
\begin{Langkah}{2}{Kumpulkan}
Tugas 5 sebelum Pertemuan 6 dimulai.
\end{Langkah}
\end{column}
\begin{column}{0.31\textwidth}
\begin{Langkah}{3}{Baca}
Modul Belajar 5 untuk mengulang, lalu bagian awal modul berikutnya.
\end{Langkah}
\end{column}
\end{columns}
\vspace{10pt}
Pertemuan 6 membahas perulangan bersarang, \texttt{break}, \texttt{continue}, dan pola bintang, sejalan dengan Algoritma Pertemuan 6.\note{Tanyakan satu hal yang paling membingungkan hari ini untuk dibuka di review minggu depan.}
\end{frame}

\SlidePenutup{Terima kasih}{Sampai jumpa di Pertemuan 6: perulangan bersarang dan pola}
\note{Slide penutup.}
\end{document}
"
\documentclass[t]{beamer}
\usetheme{ITERA}
% Mode catatan pembicara: aktif bila \catatan didefinisikan sebelum \input.
\ifdefined\catatan
  \setbeameroption{show only notes}
  \setbeamertemplate{note page}{\vspace*{20pt}\hspace*{4pt}\begin{minipage}{900pt}
    {\ITERAKickerFont\fontsize{11pt}{14pt}\selectfont\color{ITERAGoldText}CATATAN PEMBICARA \ \ |\ \ SLIDE \insertframenumber}\par\vspace{4pt}
    {\ITERADisplay\bfseries\fontsize{22pt}{28pt}\selectfont\color{ITERAStructure}\insertframetitle}\par\vspace{12pt}
    \fontsize{15pt}{23pt}\selectfont\insertnote\end{minipage}}
\fi
\tikzset{fase/.style={draw=none,fill=#1,rounded corners=3pt,minimum width=158pt,minimum height=118pt,text width=146pt,align=center}}

\renewcommand\ITERAmeeting{Pertemuan 5}
\begin{document}
\SlideJudul{IF25-11002 PRAKTIKUM PEMROGRAMAN}{Pertemuan 5\\Perulangan Bagian 1}{for, while, dan do-while dengan pola pencacah, akumulator, dan sentinel}{Tim Pengajar}{Tahun 2026. Sejajar dengan Pertemuan 5 Algoritma Pemrograman: Perulangan Bagian 1.}
\note{Buka dengan menanyakan satu hal dari pertemuan lalu yang masih membingungkan.}

\begin{frame}{Dua mata kuliah, satu minggu}
\framesubtitle{Sebelum mulai}
Topik minggu ini sama di dua mata kuliah, dengan peran yang berbeda.
\vspace{10pt}

\begin{columns}[T]
\begin{column}{0.47\textwidth}
\begin{Blok}{Algoritma: Perulangan Bagian 1}
Memodelkan pengulangan dengan flowchart dan pseudocode ULANGI dan SELAMA, lalu menelusurinya dengan trace table.
\end{Blok}
\end{column}
\begin{column}{0.47\textwidth}
\begin{BlokGelap}{Praktikum: Perulangan 1}
Menulis \texttt{for}, \texttt{while}, dan \texttt{do-while}, memakai pencacah dan akumulator, dan menghentikan perulangan dengan sentinel.
\end{BlokGelap}
\end{column}
\end{columns}
\note{Tanyakan apa yang sudah dibahas di Algoritma minggu ini. Gunakan jawabannya sebagai jembatan ke praktikum.}
\end{frame}

\begin{frame}{Saat pulang nanti, kalian bisa}
\framesubtitle{Target hari ini}
\begin{columns}[T]
\begin{column}{0.47\textwidth}
\begin{Langkah}{1}{Menerapkan}
Menulis program yang mengulang proses dengan \texttt{for}, \texttt{while}, dan \texttt{do-while}, memakai pola pencacah, akumulator, dan nilai sentinel

{\footnotesize\color{ITERAInkMuted}Sub-CPMK 5.1, CPMK0607}
\end{Langkah}
\end{column}
\begin{column}{0.47\textwidth}
\begin{Langkah}{2}{Menjelaskan}
Menelusuri perulangan dengan trace table, menentukan berapa kali badan perulangan dijalankan, dan menjelaskan penyebab off-by-one atau perulangan tak berhenti

{\footnotesize\color{ITERAInkMuted}Sub-CPMK 5.2, CPMK0608}
\end{Langkah}
\end{column}
\end{columns}
\vspace{8pt}
\begin{Catatan}
Dua kemampuan ini dinilai sama berat di Exit Ticket, tugas, UTS, dan UAS.
\end{Catatan}
\note{Bacakan target dengan pelan. Kata kuncinya menerapkan dan menjelaskan.}
\end{frame}

\begin{frame}{Ingat minggu lalu}
\framesubtitle{Review, 5 menit}
\begin{columns}[T]
\begin{column}{0.31\textwidth}
\begin{BlokGelap}{Pertanyaan 1}
Apa keluaran rantai \texttt{else if} untuk nilai batas 75?
\end{BlokGelap}
\end{column}
\begin{column}{0.31\textwidth}
\begin{BlokGelap}{Pertanyaan 2}
Mengapa perlu \texttt{break} di \texttt{switch}?
\end{BlokGelap}
\end{column}
\begin{column}{0.31\textwidth}
\begin{BlokGelap}{Pertanyaan 3}
Apa beda \texttt{if (x = 5)} dan \texttt{if (x == 5)}?
\end{BlokGelap}
\end{column}
\end{columns}
\vspace{8pt}
Jawab di kertas dulu, lalu cocokkan dengan teman sebelah.\note{Pilih dua mahasiswa untuk menjawab. Koreksi singkat, jangan mengajar ulang.}
\end{frame}

\Bagian{1}{Masalah pembuka}{Rata-rata nilai satu kelas}
\note{Masalah pembuka 10 menit.}

\begin{frame}[fragile]{Rata-rata nilai satu kelas}
\framesubtitle{Masalah minggu ini}
\begin{columns}[T]
\begin{column}{0.55\textwidth}
{Asisten praktikum harus menghitung rata-rata nilai Exit Ticket. Jumlah mahasiswa berbeda di setiap kelas, kadang 30, kadang 45.

\vspace{6pt}
Menulis \texttt{cin} sebanyak 45 kali jelas tidak masuk akal. Lagi pula banyaknya data baru diketahui saat program berjalan.\par}
\vspace{6pt}
\begin{Catatan}
\textbf{Diskusi 2 menit.} Langkah apa yang terus diulang untuk setiap nilai? Bagaimana program tahu kapan berhenti?
\end{Catatan}
\end{column}
\begin{column}{0.41\textwidth}
\begin{Keluaran}[]
Masukkan nilai (-1 untuk selesai):
70 85 60 90 -1
Banyak    : 4
Tertinggi : 90
Rata-rata : 76.25
\end{Keluaran}
\end{column}
\end{columns}
\note{Tampung beberapa jawaban tanpa menilai benar salah.}
\end{frame}

\begin{frame}{Dari langkah ke instruksi}
\framesubtitle{Sambungan dengan Algoritma}
\begin{columns}[T]
\begin{column}{0.31\textwidth}
\begin{Langkah}{1}{Temukan yang berulang}
Baca satu nilai, tambahkan ke total, tambah pencacah.
\end{Langkah}
\end{column}
\begin{column}{0.31\textwidth}
\begin{Langkah}{2}{Tentukan kapan berhenti}
Berhenti saat bertemu -1, atau setelah n kali.
\end{Langkah}
\end{column}
\begin{column}{0.31\textwidth}
\begin{Langkah}{3}{Terjemahkan}
\texttt{while} untuk syarat berhenti, \texttt{for} untuk banyak putaran yang sudah diketahui.
\end{Langkah}
\end{column}
\end{columns}
\vspace{10pt}
Langkah 1 dan 2 dilatih di Algoritma. Langkah 3 kita kerjakan hari ini dengan C++.\note{Hubungkan eksplisit ke Algoritma minggu ini.}
\end{frame}

\Bagian{2}{Coba dulu, baru dijelaskan}{25 menit eksperimen di GDB Online}
\note{Struggle 25 menit. Berkeliling, jangan menjelaskan materi dulu.}

\begin{frame}{Misi kalian}
\framesubtitle{Struggle, 25 menit}
\begin{columns}[T]
\begin{column}{0.31\textwidth}
\begin{Langkah}{1}{Hitung}
Tampilkan angka 1 sampai 10 tanpa menulis \texttt{cout} sepuluh kali.
\end{Langkah}
\end{column}
\begin{column}{0.31\textwidth}
\begin{Langkah}{2}{Jumlahkan}
Hitung 1 + 2 + ... + 100. Tebak dulu hasilnya.
\end{Langkah}
\end{column}
\begin{column}{0.31\textwidth}
\begin{Langkah}{3}{Berhenti}
Baca angka terus sampai pengguna mengetik 0.
\end{Langkah}
\end{column}
\end{columns}
\vspace{8pt}
\begin{columns}[T]
\begin{column}{0.47\textwidth}
\begin{Blok}{Boleh}
Bertanya ke teman, mengubah apa saja, sengaja membuat error.
\end{Blok}
\end{column}
\begin{column}{0.47\textwidth}
\begin{BlokAlert}{Belum dulu}
Mencari jawaban jadi di internet atau AI.
\end{BlokAlert}
\end{column}
\end{columns}
\note{Catat kesalahan yang sering muncul untuk dibahas di debrief.}
\end{frame}

\begin{frame}{Apa yang kalian temukan?}
\framesubtitle{Debrief struggle}
\begin{columns}[T]
\begin{column}{0.31\textwidth}
\begin{BlokGelap}{Pertanyaan 1}
Variabel apa yang kalian butuhkan untuk menyimpan jumlah?
\end{BlokGelap}
\end{column}
\begin{column}{0.31\textwidth}
\begin{BlokGelap}{Pertanyaan 2}
Di mana variabel itu harus diberi nilai awal?
\end{BlokGelap}
\end{column}
\begin{column}{0.31\textwidth}
\begin{BlokGelap}{Pertanyaan 3}
Apakah ada yang programnya tidak mau berhenti? Mengapa?
\end{BlokGelap}
\end{column}
\end{columns}
\vspace{10pt}
Jawaban kalian akan kita uji di bagian berikutnya.\note{Tulis jawaban di papan; buktikan atau patahkan di bagian materi.}
\end{frame}

\Bagian{3}{Mengulang dengan terkendali}{for, while, do-while, dan pola-pola perulangan}
\note{Materi 40 menit. Tunjukkan setiap contoh langsung di GDB Online.}

\begin{frame}[fragile]{Perulangan for}
\framesubtitle{Banyak putaran diketahui}
\begin{columns}[T]
\begin{column}{0.55\textwidth}
\begin{Kode}[]{for.cpp}
for (int i = 1; i <= 5; i++) {
    cout << i << " ";
}
cout << endl;
for (int i = 10; i > 0; i -= 3) {
    cout << i << " ";
}
cout << endl;
\end{Kode}
\end{column}
\begin{column}{0.41\textwidth}
\only<1>{\begin{TebakDulu}
{\ITERAKickerFont\fontsize{10pt}{12pt}\selectfont\color{ITERAAlert}TEBAK DULU}\par\vspace{4pt}
Sebelum dijalankan, tulis keluarannya di kertas.
\end{TebakDulu}
}\only<2>{\KeluaranFile[]{keluaran/P05/k001.txt}
\begin{Catatan}
Tiga bagian \texttt{for}: nilai awal; syarat lanjut; perubahan. Syarat diperiksa sebelum setiap putaran.
\end{Catatan}
}
\end{column}
\end{columns}
\note<1>{Minta mahasiswa menebak keluaran sebelum membuka slide berikutnya. Tanyakan satu atau dua tebakan yang berbeda, lalu buktikan.}
\end{frame}

\begin{frame}[fragile]{Pencacah dan akumulator}
\framesubtitle{Pola perulangan}
\begin{columns}[T]
\begin{column}{0.55\textwidth}
\begin{Kode}[]{akumulator.cpp}
int jumlah = 0, genap = 0;
for (int i = 1; i <= 10; i++) {
    jumlah += i;
    if (i % 2 == 0) {
        genap++;
    }
}
cout << "Jumlah: " << jumlah << endl;
cout << "Genap : " << genap << endl;
\end{Kode}
\end{column}
\begin{column}{0.41\textwidth}
\only<1>{\begin{TebakDulu}
{\ITERAKickerFont\fontsize{10pt}{12pt}\selectfont\color{ITERAAlert}TEBAK DULU}\par\vspace{4pt}
Sebelum dijalankan, tulis keluarannya di kertas.
\end{TebakDulu}
}\only<2>{\KeluaranFile[]{keluaran/P05/k002.txt}
\begin{Catatan}
Akumulator dan pencacah selalu diberi nilai awal sebelum perulangan, bukan di dalamnya.
\end{Catatan}
}
\end{column}
\end{columns}
\note<1>{Tanya: apa yang terjadi bila int jumlah = 0 dipindah ke dalam perulangan?}
\end{frame}

\begin{frame}[fragile]{Perulangan while}
\framesubtitle{Berhenti karena syarat}
\begin{columns}[T]
\begin{column}{0.60\textwidth}
\begin{Kode}[listing options={style=ITERACodeKecil}]{while.cpp}
int saldo = 100000, bulan = 0;
while (saldo < 500000) {
    saldo += 150000;
    bulan++;
}
cout << bulan << " bulan, saldo " << saldo << endl;
\end{Kode}
\end{column}
\begin{column}{0.36\textwidth}
\only<1>{\begin{TebakDulu}
{\ITERAKickerFont\fontsize{10pt}{12pt}\selectfont\color{ITERAAlert}TEBAK DULU}\par\vspace{4pt}
Sebelum dijalankan, tulis keluarannya di kertas.
\end{TebakDulu}
}\only<2>{\KeluaranFile[listing options={style=ITERAPlainKecil}]{keluaran/P05/k003.txt}
\begin{Catatan}
Sesuatu di dalam badan harus mengubah syarat. Bila tidak, perulangan tidak pernah berhenti.
\end{Catatan}
}
\end{column}
\end{columns}
\note<1>{Minta mahasiswa menebak keluaran sebelum membuka slide berikutnya. Tanyakan satu atau dua tebakan yang berbeda, lalu buktikan.}
\end{frame}

\begin{frame}[fragile]{Nilai sentinel}
\framesubtitle{Membaca sampai tanda berhenti}
\begin{columns}[T]
\begin{column}{0.60\textwidth}
\begin{Kode}[listing options={style=ITERACodeKecil}]{sentinel.cpp}
int nilai, banyak = 0, total = 0;
cin >> nilai;
while (nilai != -1) {
    banyak++;
    total += nilai;
    cin >> nilai;
}
cout << banyak << " data, total " << total << endl;
\end{Kode}
\end{column}
\begin{column}{0.36\textwidth}
\begin{Keluaran}[title={Masukan},listing options={style=ITERAPlainKecil}]
70 85 60 90 -1
\end{Keluaran}
\only<1>{\begin{TebakDulu}
{\ITERAKickerFont\fontsize{10pt}{12pt}\selectfont\color{ITERAAlert}TEBAK DULU}\par\vspace{4pt}
Sebelum dijalankan, tulis keluarannya di kertas.
\end{TebakDulu}
}\only<2>{\KeluaranFile[listing options={style=ITERAPlainKecil}]{keluaran/P05/k004.txt}
\begin{Catatan}
Baca sekali sebelum perulangan, lalu baca lagi di akhir badan. Dengan begitu -1 tidak ikut dihitung.
\end{Catatan}
}
\end{column}
\end{columns}
\note<1>{Minta mahasiswa menebak keluaran sebelum membuka slide berikutnya. Tanyakan satu atau dua tebakan yang berbeda, lalu buktikan.}
\end{frame}

\begin{frame}[fragile]{Perulangan do-while}
\framesubtitle{Minimal sekali jalan}
\begin{columns}[T]
\begin{column}{0.60\textwidth}
\begin{Kode}[listing options={style=ITERACodeKecil}]{dowhile.cpp}
int pin, coba = 0;
do {
    cin >> pin;
    coba++;
} while (pin != 1234 && coba < 3);
cout << (pin == 1234 ? "Masuk" : "Terkunci");
cout << " setelah " << coba << " kali" << endl;
\end{Kode}
\end{column}
\begin{column}{0.36\textwidth}
\begin{Keluaran}[title={Masukan},listing options={style=ITERAPlainKecil}]
1111 1234
\end{Keluaran}
\only<1>{\begin{TebakDulu}
{\ITERAKickerFont\fontsize{10pt}{12pt}\selectfont\color{ITERAAlert}TEBAK DULU}\par\vspace{4pt}
Sebelum dijalankan, tulis keluarannya di kertas.
\end{TebakDulu}
}\only<2>{\KeluaranFile[listing options={style=ITERAPlainKecil}]{keluaran/P05/k005.txt}
\begin{Catatan}
Syarat diperiksa sesudah badan. Jangan lupa titik koma setelah \texttt{while (...)}.
\end{Catatan}
}
\end{column}
\end{columns}
\note<1>{Minta mahasiswa menebak keluaran sebelum membuka slide berikutnya. Tanyakan satu atau dua tebakan yang berbeda, lalu buktikan.}
\end{frame}

\begin{frame}{Memilih jenis perulangan}
\framesubtitle{for, while, atau do-while}
\small
\begin{tabular}{@{}>{\raggedright\arraybackslash}p{364pt}>{\raggedright\arraybackslash}p{153pt}>{\raggedright\arraybackslash}p{345pt}@{}}
\kepala{Situasi} & \kepala{Pilihan} & \kepala{Contoh}\\
Banyak putaran diketahui & \texttt{for} & tampilkan 1 sampai n\\
\barisgenap Berhenti karena syarat, bisa nol putaran & \texttt{while} & baca sampai sentinel\\
Badan harus jalan minimal sekali & \texttt{do-while} & menu, validasi PIN\\
\garisdasar
\end{tabular}
\vspace{10pt}

Ketiganya bisa saling menggantikan. Pilih yang paling jelas menyatakan maksud kalian.
\end{frame}

\begin{frame}[fragile]{Tebak keluarannya}
\framesubtitle{Latihan menjelaskan, 3 menit}
\begin{columns}[T]
\begin{column}{0.56\textwidth}
\begin{Kode}[]{tebak.cpp}
int x = 0;
for (int i = 0; i < 4; i++) {
    x += i * 2;
}
int k = 1;
while (k < 20) {
    k *= 3;
}
cout << x << " " << k << endl;
\end{Kode}
\end{column}
\begin{column}{0.40\textwidth}
\begin{Catatan}
Tulis keluarannya di kertas, karakter demi karakter. Jangan dijalankan dulu.
\end{Catatan}
\end{column}
\end{columns}
\note{Contoh soal Bagian B. Beri waktu 3 menit sebelum slide jawaban.}
\end{frame}

\begin{frame}[fragile]{Tebak keluarannya: jawaban}
\framesubtitle{Latihan menjelaskan}
\begin{columns}[T]
\begin{column}{0.56\textwidth}
\begin{Kode}[]{tebak.cpp}
int x = 0;
for (int i = 0; i < 4; i++) {
    x += i * 2;
}
int k = 1;
while (k < 20) {
    k *= 3;
}
cout << x << " " << k << endl;
\end{Kode}
\end{column}
\begin{column}{0.40\textwidth}
\begin{Keluaran}[]
12 27
\end{Keluaran}
\begin{Catatan}
\texttt{x} = 0 + 2 + 4 + 6 = 12. \texttt{k} menjadi 3, 9, 27; berhenti saat 27 tidak lagi kurang dari 20.
\end{Catatan}
\end{column}
\end{columns}
\note{Tanyakan jawaban yang paling sering salah, lalu telusuri bersama.}
\end{frame}

\begin{frame}{Error yang sering muncul minggu ini}
\framesubtitle{Pesan diambil langsung dari g++}
\footnotesize
\begin{tabular}{@{}>{\raggedright\arraybackslash}p{208pt}>{\raggedright\arraybackslash}p{256pt}>{\raggedright\arraybackslash}p{398pt}@{}}
\kepala{Kesalahan} & \kepala{Contoh} & \kepala{Pesan pertama dari g++}\\
Koma di header for & \texttt{for (int i = 0, i < 5, i++)} & \texttt{error: expected ';' before '<' token}\\
\barisgenap Variabel di luar jangkauan & \texttt{cout << i; sesudah for} & \texttt{error: 'i' was not declared in this scope}\\
Lupa ; sesudah do-while & \texttt{\} while (x < 5)} & \texttt{error: expected ';' before 'cout'}\\
\barisgenap Perbandingan bertipe beda & \texttt{i < s.length() dengan int} & \texttt{warning: comparison of integer expressions of different signedness: 'int' and ...}\\
Badan for kosong & \texttt{for (...);} & \texttt{warning: this 'for' clause does not guard...}\\
\garisdasar
\end{tabular}
\vspace{10pt}

{\small Pesan di tabel ini dihasilkan dengan g++ dan opsi -Wall, sama seperti di cek.sh.}
\note{Minta mahasiswa memotret slide ini.}
\end{frame}

\Bagian{4}{Hands-on bertingkat}{45 menit, dari dasar sampai tantangan}
\note{Mahasiswa membuka folder 02 Hands-on/P5 - Perulangan 1 - Hands-on.}

\begin{frame}{Peta hands-on}
\framesubtitle{Folder 02 Hands-on/P5 - Perulangan 1 - Hands-on}
\small
\begin{tabular}{@{}>{\raggedright\arraybackslash}p{36pt}>{\raggedright\arraybackslash}p{267pt}>{\raggedright\arraybackslash}p{110pt}>{\raggedright\arraybackslash}p{83pt}>{\raggedright\arraybackslash}p{341pt}@{}}
\kepala{No} & \kepala{File} & \kepala{Tingkat} & \kepala{Waktu} & \kepala{Yang dilatih}\\
1 & \texttt{handson1-deret.cpp} & Dasar & 10' & \texttt{for}, akumulator\\
\barisgenap 2 & \texttt{handson2-tabungan.cpp} & Menengah & 15' & \texttt{while} dengan syarat berhenti\\
3 & \texttt{handson3-rerata.cpp} & Tantangan & 20' & trace table, perbaiki tiga kesalahan sentinel\\
\barisgenap + & \texttt{bonus-faktorial.cpp} & Pengayaan & sisa & \texttt{do-while} untuk menu, \texttt{long long}\\
\garisdasar
\end{tabular}
\vspace{10pt}

Kerjakan berurutan. Cocokkan keluaran dengan folder \texttt{expected/}; latihan yang membaca masukan memakai data di folder \texttt{input/}.\note{Saat berkeliling, minta mahasiswa yang mengerjakan latihan tantangan menjelaskan blok PENJELASAN mereka.}
\end{frame}

\begin{frame}{Cara mengecek dan aturannya}
\framesubtitle{Hands-on}
\begin{columns}[T]
\begin{column}{0.47\textwidth}
\begin{Blok}{Di GDB Online}
Salin isi file latihan, lengkapi TODO, klik Run. Bila program membaca masukan, ketik isi file di \texttt{input/}. Bandingkan dengan \texttt{expected/}.
\end{Blok}
\end{column}
\begin{column}{0.47\textwidth}
\begin{BlokGelap}{Di komputer sendiri}
Jalankan \texttt{bash cek.sh}. Skrip mengompilasi dengan \texttt{-Wall -Werror}, memberi masukan dari \texttt{input/}, lalu mencocokkan keluaran.
\end{BlokGelap}
\end{column}
\end{columns}
\vspace{6pt}
\begin{Catatan}
AI boleh untuk memahami pesan error, tidak untuk meminta solusi utuh. Exit Ticket dikerjakan tanpa bantuan apa pun.
\end{Catatan}
\note{Hands-on tidak dinilai langsung; yang dinilai Exit Ticket dan tugas.}
\end{frame}

\Bagian{5}{Exit Ticket dan tugas}{25 menit terakhir, lalu pekerjaan rumah}
\note{Mulai Exit Ticket tepat waktu.}

\begin{frame}{Exit Ticket 5}
\framesubtitle{Individual, 25 menit, tanpa AI dan internet}
\small
\begin{tabular}{@{}>{\raggedright\arraybackslash}p{60pt}>{\raggedright\arraybackslash}p{560pt}>{\raggedright\arraybackslash}p{80pt}>{\raggedright\arraybackslash}p{150pt}@{}}
\kepala{Soal} & \kepala{Isi} & \kepala{Poin} & \kepala{CPMK}\\
A1 & Tulis program yang menampilkan kelipatan k dari 1 sampai n beserta banyaknya & 25 & CPMK0607\\
\barisgenap A2 & Tulis program yang membaca bilangan sampai 0 lalu menampilkan jumlah bilangan positif & 25 & CPMK0607\\
B1 & Buat trace table sebuah perulangan dan tentukan keluarannya & 20 & CPMK0608\\
\barisgenap B2 & Jelaskan mengapa sebuah perulangan tidak berhenti dan cara memperbaikinya & 20 & CPMK0608\\
B3 & Tentukan berapa kali badan perulangan dijalankan dan jelaskan alasannya & 10 & CPMK0608\\
\garisdasar
\end{tabular}
\vspace{10pt}

{\small Soal A dinilai dari keluaran yang tepat sama. Soal B dinilai dari ketepatan dan alasan. Pelanggaran aturan membuat nilai Exit Ticket ini 0.}
\note{Soal lengkap dibagikan terpisah dan tidak disimpan di repo publik.}
\end{frame}

\begin{frame}{Tugas 5: Program Statistik Nilai}
\framesubtitle{Dikumpulkan sebelum Pertemuan 6}
\begin{columns}[T]
\begin{column}{0.47\textwidth}
\begin{Blok}{Bagian A, program, 50 poin}
Buat program statistik nilai dengan ketentuan berikut. Ketentuannya di slide berikut.
\end{Blok}
\end{column}
\begin{column}{0.47\textwidth}
\begin{BlokGelap}{Bagian B, penjelasan, 50 poin}
Komentar maksimal 150 kata dan trace table untuk n = 3 dengan nilai 80, 70, 90. Jelaskan alasan memilih \texttt{for}, \texttt{while}, atau \texttt{do-while} di setiap bagian, dan satu kesalahan off-by-one atau perulangan tak berhenti yang kalian temui.
\end{BlokGelap}
\end{column}
\end{columns}
\vspace{8pt}
{\small Data di tugas adalah data kalian sendiri. Rubrik lengkap ada di Modul Belajar 5.}
\note{Ingatkan bahwa penjelasan disalin dari AI mudah dikenali karena tidak menyebut pengalaman mereka sendiri.}
\end{frame}

\begin{frame}{Ketentuan Tugas 5}
\framesubtitle{Program Statistik Nilai}
\small
\begin{tabular}{@{}>{\raggedright\arraybackslash}p{87pt}>{\raggedright\arraybackslash}p{788pt}@{}}
\kepala{No} & \kepala{Ketentuan Bagian A}\\
1 & Baca jumlah siswa n, ulangi sampai n > 0\\
\barisgenap 2 & Baca nilai n siswa; setiap nilai harus 0 sampai 100, ulangi bila tidak valid\\
3 & Tampilkan nilai tertinggi, nilai terendah, dan rata-rata kelas\\
\barisgenap 4 & Tampilkan jumlah siswa LULUS (nilai 75 ke atas) dan TIDAK LULUS\\
5 & Menu \texttt{do-while} untuk mengulang atau keluar\\
\barisgenap Bonus & Tampilkan persentase kelulusan dengan dua desimal.\\
\garisdasar
\end{tabular}
\note{Bacakan ketentuan satu per satu dan tanyakan apakah ada yang belum jelas.}
\end{frame}

\begin{frame}{Rubrik Tugas 5}
\framesubtitle{Empat level, sama di semua instrumen}
\footnotesize
\begin{tabular}{@{}>{\raggedright\arraybackslash}p{166pt}>{\raggedright\arraybackslash}p{44pt}>{\raggedright\arraybackslash}p{166pt}>{\raggedright\arraybackslash}p{156pt}>{\raggedright\arraybackslash}p{146pt}>{\raggedright\arraybackslash}p{146pt}@{}}
\kepala{Kriteria} & \kepala{Poin} & \kepala{Sangat baik 85--100} & \kepala{Baik 70--84} & \kepala{Cukup 55--69} & \kepala{Kurang <55}\\
A1 Program berjalan & 20 & tanpa error dan peringatan & ada peringatan & perlu satu perbaikan & tidak terkompilasi\\
\barisgenap A2 Validasi dan perulangan & 15 & validasi n dan nilai, menu berulang & satu validasi kurang & dua kurang & tanpa validasi\\
A3 Ketepatan statistik & 15 & semua tepat termasuk n = 1 & satu salah & dua salah & sebagian besar salah\\
\barisgenap B1 Trace table dan alasan & 30 & trace lengkap dan alasan tepat & satu baris trace salah & trace tidak lengkap & tidak ada atau disalin\\
B2 Cerita error dan pengujian & 20 & pesan, sebab, perbaikan, uji & pesan tidak dikutip & hanya perbaikan & tidak ada\\
\garisdasar
\end{tabular}
\vspace{8pt}

{\small Nilai kriteria = poin × skor level. Bagian A total 50, Bagian B total 50.}
\note{Rubrik ini sama strukturnya setiap minggu; yang berubah hanya isi kriteria A2, A3, dan B1.}
\end{frame}

\begin{frame}{Sebelum Pertemuan 6}
\framesubtitle{Persiapan}
\begin{columns}[T]
\begin{column}{0.31\textwidth}
\begin{Langkah}{1}{Selesaikan}
Hands-on yang belum lulus \texttt{cek.sh}, terutama latihan tantangan.
\end{Langkah}
\end{column}
\begin{column}{0.31\textwidth}
\begin{Langkah}{2}{Kumpulkan}
Tugas 5 sebelum Pertemuan 6 dimulai.
\end{Langkah}
\end{column}
\begin{column}{0.31\textwidth}
\begin{Langkah}{3}{Baca}
Modul Belajar 5 untuk mengulang, lalu bagian awal modul berikutnya.
\end{Langkah}
\end{column}
\end{columns}
\vspace{10pt}
Pertemuan 6 membahas perulangan bersarang, \texttt{break}, \texttt{continue}, dan pola bintang, sejalan dengan Algoritma Pertemuan 6.\note{Tanyakan satu hal yang paling membingungkan hari ini untuk dibuka di review minggu depan.}
\end{frame}

\SlidePenutup{Terima kasih}{Sampai jumpa di Pertemuan 6: perulangan bersarang dan pola}
\note{Slide penutup.}
\end{document}
"
\documentclass[t]{beamer}
\usetheme{ITERA}
% Mode catatan pembicara: aktif bila \catatan didefinisikan sebelum \input.
\ifdefined\catatan
  \setbeameroption{show only notes}
  \setbeamertemplate{note page}{\vspace*{20pt}\hspace*{4pt}\begin{minipage}{900pt}
    {\ITERAKickerFont\fontsize{11pt}{14pt}\selectfont\color{ITERAGoldText}CATATAN PEMBICARA \ \ |\ \ SLIDE \insertframenumber}\par\vspace{4pt}
    {\ITERADisplay\bfseries\fontsize{22pt}{28pt}\selectfont\color{ITERAStructure}\insertframetitle}\par\vspace{12pt}
    \fontsize{15pt}{23pt}\selectfont\insertnote\end{minipage}}
\fi
\tikzset{fase/.style={draw=none,fill=#1,rounded corners=3pt,minimum width=158pt,minimum height=118pt,text width=146pt,align=center}}

\renewcommand\ITERAmeeting{Pertemuan 5}
\begin{document}
\SlideJudul{IF25-11002 PRAKTIKUM PEMROGRAMAN}{Pertemuan 5\\Perulangan Bagian 1}{for, while, dan do-while dengan pola pencacah, akumulator, dan sentinel}{Tim Pengajar}{Tahun 2026. Sejajar dengan Pertemuan 5 Algoritma Pemrograman: Perulangan Bagian 1.}
\note{Buka dengan menanyakan satu hal dari pertemuan lalu yang masih membingungkan.}

\begin{frame}{Dua mata kuliah, satu minggu}
\framesubtitle{Sebelum mulai}
Topik minggu ini sama di dua mata kuliah, dengan peran yang berbeda.
\vspace{10pt}

\begin{columns}[T]
\begin{column}{0.47\textwidth}
\begin{Blok}{Algoritma: Perulangan Bagian 1}
Memodelkan pengulangan dengan flowchart dan pseudocode ULANGI dan SELAMA, lalu menelusurinya dengan trace table.
\end{Blok}
\end{column}
\begin{column}{0.47\textwidth}
\begin{BlokGelap}{Praktikum: Perulangan 1}
Menulis \texttt{for}, \texttt{while}, dan \texttt{do-while}, memakai pencacah dan akumulator, dan menghentikan perulangan dengan sentinel.
\end{BlokGelap}
\end{column}
\end{columns}
\note{Tanyakan apa yang sudah dibahas di Algoritma minggu ini. Gunakan jawabannya sebagai jembatan ke praktikum.}
\end{frame}

\begin{frame}{Saat pulang nanti, kalian bisa}
\framesubtitle{Target hari ini}
\begin{columns}[T]
\begin{column}{0.47\textwidth}
\begin{Langkah}{1}{Menerapkan}
Menulis program yang mengulang proses dengan \texttt{for}, \texttt{while}, dan \texttt{do-while}, memakai pola pencacah, akumulator, dan nilai sentinel

{\footnotesize\color{ITERAInkMuted}Sub-CPMK 5.1, CPMK0607}
\end{Langkah}
\end{column}
\begin{column}{0.47\textwidth}
\begin{Langkah}{2}{Menjelaskan}
Menelusuri perulangan dengan trace table, menentukan berapa kali badan perulangan dijalankan, dan menjelaskan penyebab off-by-one atau perulangan tak berhenti

{\footnotesize\color{ITERAInkMuted}Sub-CPMK 5.2, CPMK0608}
\end{Langkah}
\end{column}
\end{columns}
\vspace{8pt}
\begin{Catatan}
Dua kemampuan ini dinilai sama berat di Exit Ticket, tugas, UTS, dan UAS.
\end{Catatan}
\note{Bacakan target dengan pelan. Kata kuncinya menerapkan dan menjelaskan.}
\end{frame}

\begin{frame}{Ingat minggu lalu}
\framesubtitle{Review, 5 menit}
\begin{columns}[T]
\begin{column}{0.31\textwidth}
\begin{BlokGelap}{Pertanyaan 1}
Apa keluaran rantai \texttt{else if} untuk nilai batas 75?
\end{BlokGelap}
\end{column}
\begin{column}{0.31\textwidth}
\begin{BlokGelap}{Pertanyaan 2}
Mengapa perlu \texttt{break} di \texttt{switch}?
\end{BlokGelap}
\end{column}
\begin{column}{0.31\textwidth}
\begin{BlokGelap}{Pertanyaan 3}
Apa beda \texttt{if (x = 5)} dan \texttt{if (x == 5)}?
\end{BlokGelap}
\end{column}
\end{columns}
\vspace{8pt}
Jawab di kertas dulu, lalu cocokkan dengan teman sebelah.\note{Pilih dua mahasiswa untuk menjawab. Koreksi singkat, jangan mengajar ulang.}
\end{frame}

\Bagian{1}{Masalah pembuka}{Rata-rata nilai satu kelas}
\note{Masalah pembuka 10 menit.}

\begin{frame}[fragile]{Rata-rata nilai satu kelas}
\framesubtitle{Masalah minggu ini}
\begin{columns}[T]
\begin{column}{0.55\textwidth}
{Asisten praktikum harus menghitung rata-rata nilai Exit Ticket. Jumlah mahasiswa berbeda di setiap kelas, kadang 30, kadang 45.

\vspace{6pt}
Menulis \texttt{cin} sebanyak 45 kali jelas tidak masuk akal. Lagi pula banyaknya data baru diketahui saat program berjalan.\par}
\vspace{6pt}
\begin{Catatan}
\textbf{Diskusi 2 menit.} Langkah apa yang terus diulang untuk setiap nilai? Bagaimana program tahu kapan berhenti?
\end{Catatan}
\end{column}
\begin{column}{0.41\textwidth}
\begin{Keluaran}[]
Masukkan nilai (-1 untuk selesai):
70 85 60 90 -1
Banyak    : 4
Tertinggi : 90
Rata-rata : 76.25
\end{Keluaran}
\end{column}
\end{columns}
\note{Tampung beberapa jawaban tanpa menilai benar salah.}
\end{frame}

\begin{frame}{Dari langkah ke instruksi}
\framesubtitle{Sambungan dengan Algoritma}
\begin{columns}[T]
\begin{column}{0.31\textwidth}
\begin{Langkah}{1}{Temukan yang berulang}
Baca satu nilai, tambahkan ke total, tambah pencacah.
\end{Langkah}
\end{column}
\begin{column}{0.31\textwidth}
\begin{Langkah}{2}{Tentukan kapan berhenti}
Berhenti saat bertemu -1, atau setelah n kali.
\end{Langkah}
\end{column}
\begin{column}{0.31\textwidth}
\begin{Langkah}{3}{Terjemahkan}
\texttt{while} untuk syarat berhenti, \texttt{for} untuk banyak putaran yang sudah diketahui.
\end{Langkah}
\end{column}
\end{columns}
\vspace{10pt}
Langkah 1 dan 2 dilatih di Algoritma. Langkah 3 kita kerjakan hari ini dengan C++.\note{Hubungkan eksplisit ke Algoritma minggu ini.}
\end{frame}

\Bagian{2}{Coba dulu, baru dijelaskan}{25 menit eksperimen di GDB Online}
\note{Struggle 25 menit. Berkeliling, jangan menjelaskan materi dulu.}

\begin{frame}{Misi kalian}
\framesubtitle{Struggle, 25 menit}
\begin{columns}[T]
\begin{column}{0.31\textwidth}
\begin{Langkah}{1}{Hitung}
Tampilkan angka 1 sampai 10 tanpa menulis \texttt{cout} sepuluh kali.
\end{Langkah}
\end{column}
\begin{column}{0.31\textwidth}
\begin{Langkah}{2}{Jumlahkan}
Hitung 1 + 2 + ... + 100. Tebak dulu hasilnya.
\end{Langkah}
\end{column}
\begin{column}{0.31\textwidth}
\begin{Langkah}{3}{Berhenti}
Baca angka terus sampai pengguna mengetik 0.
\end{Langkah}
\end{column}
\end{columns}
\vspace{8pt}
\begin{columns}[T]
\begin{column}{0.47\textwidth}
\begin{Blok}{Boleh}
Bertanya ke teman, mengubah apa saja, sengaja membuat error.
\end{Blok}
\end{column}
\begin{column}{0.47\textwidth}
\begin{BlokAlert}{Belum dulu}
Mencari jawaban jadi di internet atau AI.
\end{BlokAlert}
\end{column}
\end{columns}
\note{Catat kesalahan yang sering muncul untuk dibahas di debrief.}
\end{frame}

\begin{frame}{Apa yang kalian temukan?}
\framesubtitle{Debrief struggle}
\begin{columns}[T]
\begin{column}{0.31\textwidth}
\begin{BlokGelap}{Pertanyaan 1}
Variabel apa yang kalian butuhkan untuk menyimpan jumlah?
\end{BlokGelap}
\end{column}
\begin{column}{0.31\textwidth}
\begin{BlokGelap}{Pertanyaan 2}
Di mana variabel itu harus diberi nilai awal?
\end{BlokGelap}
\end{column}
\begin{column}{0.31\textwidth}
\begin{BlokGelap}{Pertanyaan 3}
Apakah ada yang programnya tidak mau berhenti? Mengapa?
\end{BlokGelap}
\end{column}
\end{columns}
\vspace{10pt}
Jawaban kalian akan kita uji di bagian berikutnya.\note{Tulis jawaban di papan; buktikan atau patahkan di bagian materi.}
\end{frame}

\Bagian{3}{Mengulang dengan terkendali}{for, while, do-while, dan pola-pola perulangan}
\note{Materi 40 menit. Tunjukkan setiap contoh langsung di GDB Online.}

\begin{frame}[fragile]{Perulangan for}
\framesubtitle{Banyak putaran diketahui}
\begin{columns}[T]
\begin{column}{0.55\textwidth}
\begin{Kode}[]{for.cpp}
for (int i = 1; i <= 5; i++) {
    cout << i << " ";
}
cout << endl;
for (int i = 10; i > 0; i -= 3) {
    cout << i << " ";
}
cout << endl;
\end{Kode}
\end{column}
\begin{column}{0.41\textwidth}
\only<1>{\begin{TebakDulu}
{\ITERAKickerFont\fontsize{10pt}{12pt}\selectfont\color{ITERAAlert}TEBAK DULU}\par\vspace{4pt}
Sebelum dijalankan, tulis keluarannya di kertas.
\end{TebakDulu}
}\only<2>{\KeluaranFile[]{keluaran/P05/k001.txt}
\begin{Catatan}
Tiga bagian \texttt{for}: nilai awal; syarat lanjut; perubahan. Syarat diperiksa sebelum setiap putaran.
\end{Catatan}
}
\end{column}
\end{columns}
\note<1>{Minta mahasiswa menebak keluaran sebelum membuka slide berikutnya. Tanyakan satu atau dua tebakan yang berbeda, lalu buktikan.}
\end{frame}

\begin{frame}[fragile]{Pencacah dan akumulator}
\framesubtitle{Pola perulangan}
\begin{columns}[T]
\begin{column}{0.55\textwidth}
\begin{Kode}[]{akumulator.cpp}
int jumlah = 0, genap = 0;
for (int i = 1; i <= 10; i++) {
    jumlah += i;
    if (i % 2 == 0) {
        genap++;
    }
}
cout << "Jumlah: " << jumlah << endl;
cout << "Genap : " << genap << endl;
\end{Kode}
\end{column}
\begin{column}{0.41\textwidth}
\only<1>{\begin{TebakDulu}
{\ITERAKickerFont\fontsize{10pt}{12pt}\selectfont\color{ITERAAlert}TEBAK DULU}\par\vspace{4pt}
Sebelum dijalankan, tulis keluarannya di kertas.
\end{TebakDulu}
}\only<2>{\KeluaranFile[]{keluaran/P05/k002.txt}
\begin{Catatan}
Akumulator dan pencacah selalu diberi nilai awal sebelum perulangan, bukan di dalamnya.
\end{Catatan}
}
\end{column}
\end{columns}
\note<1>{Tanya: apa yang terjadi bila int jumlah = 0 dipindah ke dalam perulangan?}
\end{frame}

\begin{frame}[fragile]{Perulangan while}
\framesubtitle{Berhenti karena syarat}
\begin{columns}[T]
\begin{column}{0.60\textwidth}
\begin{Kode}[listing options={style=ITERACodeKecil}]{while.cpp}
int saldo = 100000, bulan = 0;
while (saldo < 500000) {
    saldo += 150000;
    bulan++;
}
cout << bulan << " bulan, saldo " << saldo << endl;
\end{Kode}
\end{column}
\begin{column}{0.36\textwidth}
\only<1>{\begin{TebakDulu}
{\ITERAKickerFont\fontsize{10pt}{12pt}\selectfont\color{ITERAAlert}TEBAK DULU}\par\vspace{4pt}
Sebelum dijalankan, tulis keluarannya di kertas.
\end{TebakDulu}
}\only<2>{\KeluaranFile[listing options={style=ITERAPlainKecil}]{keluaran/P05/k003.txt}
\begin{Catatan}
Sesuatu di dalam badan harus mengubah syarat. Bila tidak, perulangan tidak pernah berhenti.
\end{Catatan}
}
\end{column}
\end{columns}
\note<1>{Minta mahasiswa menebak keluaran sebelum membuka slide berikutnya. Tanyakan satu atau dua tebakan yang berbeda, lalu buktikan.}
\end{frame}

\begin{frame}[fragile]{Nilai sentinel}
\framesubtitle{Membaca sampai tanda berhenti}
\begin{columns}[T]
\begin{column}{0.60\textwidth}
\begin{Kode}[listing options={style=ITERACodeKecil}]{sentinel.cpp}
int nilai, banyak = 0, total = 0;
cin >> nilai;
while (nilai != -1) {
    banyak++;
    total += nilai;
    cin >> nilai;
}
cout << banyak << " data, total " << total << endl;
\end{Kode}
\end{column}
\begin{column}{0.36\textwidth}
\begin{Keluaran}[title={Masukan},listing options={style=ITERAPlainKecil}]
70 85 60 90 -1
\end{Keluaran}
\only<1>{\begin{TebakDulu}
{\ITERAKickerFont\fontsize{10pt}{12pt}\selectfont\color{ITERAAlert}TEBAK DULU}\par\vspace{4pt}
Sebelum dijalankan, tulis keluarannya di kertas.
\end{TebakDulu}
}\only<2>{\KeluaranFile[listing options={style=ITERAPlainKecil}]{keluaran/P05/k004.txt}
\begin{Catatan}
Baca sekali sebelum perulangan, lalu baca lagi di akhir badan. Dengan begitu -1 tidak ikut dihitung.
\end{Catatan}
}
\end{column}
\end{columns}
\note<1>{Minta mahasiswa menebak keluaran sebelum membuka slide berikutnya. Tanyakan satu atau dua tebakan yang berbeda, lalu buktikan.}
\end{frame}

\begin{frame}[fragile]{Perulangan do-while}
\framesubtitle{Minimal sekali jalan}
\begin{columns}[T]
\begin{column}{0.60\textwidth}
\begin{Kode}[listing options={style=ITERACodeKecil}]{dowhile.cpp}
int pin, coba = 0;
do {
    cin >> pin;
    coba++;
} while (pin != 1234 && coba < 3);
cout << (pin == 1234 ? "Masuk" : "Terkunci");
cout << " setelah " << coba << " kali" << endl;
\end{Kode}
\end{column}
\begin{column}{0.36\textwidth}
\begin{Keluaran}[title={Masukan},listing options={style=ITERAPlainKecil}]
1111 1234
\end{Keluaran}
\only<1>{\begin{TebakDulu}
{\ITERAKickerFont\fontsize{10pt}{12pt}\selectfont\color{ITERAAlert}TEBAK DULU}\par\vspace{4pt}
Sebelum dijalankan, tulis keluarannya di kertas.
\end{TebakDulu}
}\only<2>{\KeluaranFile[listing options={style=ITERAPlainKecil}]{keluaran/P05/k005.txt}
\begin{Catatan}
Syarat diperiksa sesudah badan. Jangan lupa titik koma setelah \texttt{while (...)}.
\end{Catatan}
}
\end{column}
\end{columns}
\note<1>{Minta mahasiswa menebak keluaran sebelum membuka slide berikutnya. Tanyakan satu atau dua tebakan yang berbeda, lalu buktikan.}
\end{frame}

\begin{frame}{Memilih jenis perulangan}
\framesubtitle{for, while, atau do-while}
\small
\begin{tabular}{@{}>{\raggedright\arraybackslash}p{364pt}>{\raggedright\arraybackslash}p{153pt}>{\raggedright\arraybackslash}p{345pt}@{}}
\kepala{Situasi} & \kepala{Pilihan} & \kepala{Contoh}\\
Banyak putaran diketahui & \texttt{for} & tampilkan 1 sampai n\\
\barisgenap Berhenti karena syarat, bisa nol putaran & \texttt{while} & baca sampai sentinel\\
Badan harus jalan minimal sekali & \texttt{do-while} & menu, validasi PIN\\
\garisdasar
\end{tabular}
\vspace{10pt}

Ketiganya bisa saling menggantikan. Pilih yang paling jelas menyatakan maksud kalian.
\end{frame}

\begin{frame}[fragile]{Tebak keluarannya}
\framesubtitle{Latihan menjelaskan, 3 menit}
\begin{columns}[T]
\begin{column}{0.56\textwidth}
\begin{Kode}[]{tebak.cpp}
int x = 0;
for (int i = 0; i < 4; i++) {
    x += i * 2;
}
int k = 1;
while (k < 20) {
    k *= 3;
}
cout << x << " " << k << endl;
\end{Kode}
\end{column}
\begin{column}{0.40\textwidth}
\begin{Catatan}
Tulis keluarannya di kertas, karakter demi karakter. Jangan dijalankan dulu.
\end{Catatan}
\end{column}
\end{columns}
\note{Contoh soal Bagian B. Beri waktu 3 menit sebelum slide jawaban.}
\end{frame}

\begin{frame}[fragile]{Tebak keluarannya: jawaban}
\framesubtitle{Latihan menjelaskan}
\begin{columns}[T]
\begin{column}{0.56\textwidth}
\begin{Kode}[]{tebak.cpp}
int x = 0;
for (int i = 0; i < 4; i++) {
    x += i * 2;
}
int k = 1;
while (k < 20) {
    k *= 3;
}
cout << x << " " << k << endl;
\end{Kode}
\end{column}
\begin{column}{0.40\textwidth}
\begin{Keluaran}[]
12 27
\end{Keluaran}
\begin{Catatan}
\texttt{x} = 0 + 2 + 4 + 6 = 12. \texttt{k} menjadi 3, 9, 27; berhenti saat 27 tidak lagi kurang dari 20.
\end{Catatan}
\end{column}
\end{columns}
\note{Tanyakan jawaban yang paling sering salah, lalu telusuri bersama.}
\end{frame}

\begin{frame}{Error yang sering muncul minggu ini}
\framesubtitle{Pesan diambil langsung dari g++}
\footnotesize
\begin{tabular}{@{}>{\raggedright\arraybackslash}p{208pt}>{\raggedright\arraybackslash}p{256pt}>{\raggedright\arraybackslash}p{398pt}@{}}
\kepala{Kesalahan} & \kepala{Contoh} & \kepala{Pesan pertama dari g++}\\
Koma di header for & \texttt{for (int i = 0, i < 5, i++)} & \texttt{error: expected ';' before '<' token}\\
\barisgenap Variabel di luar jangkauan & \texttt{cout << i; sesudah for} & \texttt{error: 'i' was not declared in this scope}\\
Lupa ; sesudah do-while & \texttt{\} while (x < 5)} & \texttt{error: expected ';' before 'cout'}\\
\barisgenap Perbandingan bertipe beda & \texttt{i < s.length() dengan int} & \texttt{warning: comparison of integer expressions of different signedness: 'int' and ...}\\
Badan for kosong & \texttt{for (...);} & \texttt{warning: this 'for' clause does not guard...}\\
\garisdasar
\end{tabular}
\vspace{10pt}

{\small Pesan di tabel ini dihasilkan dengan g++ dan opsi -Wall, sama seperti di cek.sh.}
\note{Minta mahasiswa memotret slide ini.}
\end{frame}

\Bagian{4}{Hands-on bertingkat}{45 menit, dari dasar sampai tantangan}
\note{Mahasiswa membuka folder 02 Hands-on/P5 - Perulangan 1 - Hands-on.}

\begin{frame}{Peta hands-on}
\framesubtitle{Folder 02 Hands-on/P5 - Perulangan 1 - Hands-on}
\small
\begin{tabular}{@{}>{\raggedright\arraybackslash}p{36pt}>{\raggedright\arraybackslash}p{267pt}>{\raggedright\arraybackslash}p{110pt}>{\raggedright\arraybackslash}p{83pt}>{\raggedright\arraybackslash}p{341pt}@{}}
\kepala{No} & \kepala{File} & \kepala{Tingkat} & \kepala{Waktu} & \kepala{Yang dilatih}\\
1 & \texttt{handson1-deret.cpp} & Dasar & 10' & \texttt{for}, akumulator\\
\barisgenap 2 & \texttt{handson2-tabungan.cpp} & Menengah & 15' & \texttt{while} dengan syarat berhenti\\
3 & \texttt{handson3-rerata.cpp} & Tantangan & 20' & trace table, perbaiki tiga kesalahan sentinel\\
\barisgenap + & \texttt{bonus-faktorial.cpp} & Pengayaan & sisa & \texttt{do-while} untuk menu, \texttt{long long}\\
\garisdasar
\end{tabular}
\vspace{10pt}

Kerjakan berurutan. Cocokkan keluaran dengan folder \texttt{expected/}; latihan yang membaca masukan memakai data di folder \texttt{input/}.\note{Saat berkeliling, minta mahasiswa yang mengerjakan latihan tantangan menjelaskan blok PENJELASAN mereka.}
\end{frame}

\begin{frame}{Cara mengecek dan aturannya}
\framesubtitle{Hands-on}
\begin{columns}[T]
\begin{column}{0.47\textwidth}
\begin{Blok}{Di GDB Online}
Salin isi file latihan, lengkapi TODO, klik Run. Bila program membaca masukan, ketik isi file di \texttt{input/}. Bandingkan dengan \texttt{expected/}.
\end{Blok}
\end{column}
\begin{column}{0.47\textwidth}
\begin{BlokGelap}{Di komputer sendiri}
Jalankan \texttt{bash cek.sh}. Skrip mengompilasi dengan \texttt{-Wall -Werror}, memberi masukan dari \texttt{input/}, lalu mencocokkan keluaran.
\end{BlokGelap}
\end{column}
\end{columns}
\vspace{6pt}
\begin{Catatan}
AI boleh untuk memahami pesan error, tidak untuk meminta solusi utuh. Exit Ticket dikerjakan tanpa bantuan apa pun.
\end{Catatan}
\note{Hands-on tidak dinilai langsung; yang dinilai Exit Ticket dan tugas.}
\end{frame}

\Bagian{5}{Exit Ticket dan tugas}{25 menit terakhir, lalu pekerjaan rumah}
\note{Mulai Exit Ticket tepat waktu.}

\begin{frame}{Exit Ticket 5}
\framesubtitle{Individual, 25 menit, tanpa AI dan internet}
\small
\begin{tabular}{@{}>{\raggedright\arraybackslash}p{60pt}>{\raggedright\arraybackslash}p{560pt}>{\raggedright\arraybackslash}p{80pt}>{\raggedright\arraybackslash}p{150pt}@{}}
\kepala{Soal} & \kepala{Isi} & \kepala{Poin} & \kepala{CPMK}\\
A1 & Tulis program yang menampilkan kelipatan k dari 1 sampai n beserta banyaknya & 25 & CPMK0607\\
\barisgenap A2 & Tulis program yang membaca bilangan sampai 0 lalu menampilkan jumlah bilangan positif & 25 & CPMK0607\\
B1 & Buat trace table sebuah perulangan dan tentukan keluarannya & 20 & CPMK0608\\
\barisgenap B2 & Jelaskan mengapa sebuah perulangan tidak berhenti dan cara memperbaikinya & 20 & CPMK0608\\
B3 & Tentukan berapa kali badan perulangan dijalankan dan jelaskan alasannya & 10 & CPMK0608\\
\garisdasar
\end{tabular}
\vspace{10pt}

{\small Soal A dinilai dari keluaran yang tepat sama. Soal B dinilai dari ketepatan dan alasan. Pelanggaran aturan membuat nilai Exit Ticket ini 0.}
\note{Soal lengkap dibagikan terpisah dan tidak disimpan di repo publik.}
\end{frame}

\begin{frame}{Tugas 5: Program Statistik Nilai}
\framesubtitle{Dikumpulkan sebelum Pertemuan 6}
\begin{columns}[T]
\begin{column}{0.47\textwidth}
\begin{Blok}{Bagian A, program, 50 poin}
Buat program statistik nilai dengan ketentuan berikut. Ketentuannya di slide berikut.
\end{Blok}
\end{column}
\begin{column}{0.47\textwidth}
\begin{BlokGelap}{Bagian B, penjelasan, 50 poin}
Komentar maksimal 150 kata dan trace table untuk n = 3 dengan nilai 80, 70, 90. Jelaskan alasan memilih \texttt{for}, \texttt{while}, atau \texttt{do-while} di setiap bagian, dan satu kesalahan off-by-one atau perulangan tak berhenti yang kalian temui.
\end{BlokGelap}
\end{column}
\end{columns}
\vspace{8pt}
{\small Data di tugas adalah data kalian sendiri. Rubrik lengkap ada di Modul Belajar 5.}
\note{Ingatkan bahwa penjelasan disalin dari AI mudah dikenali karena tidak menyebut pengalaman mereka sendiri.}
\end{frame}

\begin{frame}{Ketentuan Tugas 5}
\framesubtitle{Program Statistik Nilai}
\small
\begin{tabular}{@{}>{\raggedright\arraybackslash}p{87pt}>{\raggedright\arraybackslash}p{788pt}@{}}
\kepala{No} & \kepala{Ketentuan Bagian A}\\
1 & Baca jumlah siswa n, ulangi sampai n > 0\\
\barisgenap 2 & Baca nilai n siswa; setiap nilai harus 0 sampai 100, ulangi bila tidak valid\\
3 & Tampilkan nilai tertinggi, nilai terendah, dan rata-rata kelas\\
\barisgenap 4 & Tampilkan jumlah siswa LULUS (nilai 75 ke atas) dan TIDAK LULUS\\
5 & Menu \texttt{do-while} untuk mengulang atau keluar\\
\barisgenap Bonus & Tampilkan persentase kelulusan dengan dua desimal.\\
\garisdasar
\end{tabular}
\note{Bacakan ketentuan satu per satu dan tanyakan apakah ada yang belum jelas.}
\end{frame}

\begin{frame}{Rubrik Tugas 5}
\framesubtitle{Empat level, sama di semua instrumen}
\footnotesize
\begin{tabular}{@{}>{\raggedright\arraybackslash}p{166pt}>{\raggedright\arraybackslash}p{44pt}>{\raggedright\arraybackslash}p{166pt}>{\raggedright\arraybackslash}p{156pt}>{\raggedright\arraybackslash}p{146pt}>{\raggedright\arraybackslash}p{146pt}@{}}
\kepala{Kriteria} & \kepala{Poin} & \kepala{Sangat baik 85--100} & \kepala{Baik 70--84} & \kepala{Cukup 55--69} & \kepala{Kurang <55}\\
A1 Program berjalan & 20 & tanpa error dan peringatan & ada peringatan & perlu satu perbaikan & tidak terkompilasi\\
\barisgenap A2 Validasi dan perulangan & 15 & validasi n dan nilai, menu berulang & satu validasi kurang & dua kurang & tanpa validasi\\
A3 Ketepatan statistik & 15 & semua tepat termasuk n = 1 & satu salah & dua salah & sebagian besar salah\\
\barisgenap B1 Trace table dan alasan & 30 & trace lengkap dan alasan tepat & satu baris trace salah & trace tidak lengkap & tidak ada atau disalin\\
B2 Cerita error dan pengujian & 20 & pesan, sebab, perbaikan, uji & pesan tidak dikutip & hanya perbaikan & tidak ada\\
\garisdasar
\end{tabular}
\vspace{8pt}

{\small Nilai kriteria = poin × skor level. Bagian A total 50, Bagian B total 50.}
\note{Rubrik ini sama strukturnya setiap minggu; yang berubah hanya isi kriteria A2, A3, dan B1.}
\end{frame}

\begin{frame}{Sebelum Pertemuan 6}
\framesubtitle{Persiapan}
\begin{columns}[T]
\begin{column}{0.31\textwidth}
\begin{Langkah}{1}{Selesaikan}
Hands-on yang belum lulus \texttt{cek.sh}, terutama latihan tantangan.
\end{Langkah}
\end{column}
\begin{column}{0.31\textwidth}
\begin{Langkah}{2}{Kumpulkan}
Tugas 5 sebelum Pertemuan 6 dimulai.
\end{Langkah}
\end{column}
\begin{column}{0.31\textwidth}
\begin{Langkah}{3}{Baca}
Modul Belajar 5 untuk mengulang, lalu bagian awal modul berikutnya.
\end{Langkah}
\end{column}
\end{columns}
\vspace{10pt}
Pertemuan 6 membahas perulangan bersarang, \texttt{break}, \texttt{continue}, dan pola bintang, sejalan dengan Algoritma Pertemuan 6.\note{Tanyakan satu hal yang paling membingungkan hari ini untuk dibuka di review minggu depan.}
\end{frame}

\SlidePenutup{Terima kasih}{Sampai jumpa di Pertemuan 6: perulangan bersarang dan pola}
\note{Slide penutup.}
\end{document}
"
\documentclass[t]{beamer}
\usetheme{ITERA}
% Mode catatan pembicara: aktif bila \catatan didefinisikan sebelum \input.
\ifdefined\catatan
  \setbeameroption{show only notes}
  \setbeamertemplate{note page}{\vspace*{20pt}\hspace*{4pt}\begin{minipage}{900pt}
    {\ITERAKickerFont\fontsize{11pt}{14pt}\selectfont\color{ITERAGoldText}CATATAN PEMBICARA \ \ |\ \ SLIDE \insertframenumber}\par\vspace{4pt}
    {\ITERADisplay\bfseries\fontsize{22pt}{28pt}\selectfont\color{ITERAStructure}\insertframetitle}\par\vspace{12pt}
    \fontsize{15pt}{23pt}\selectfont\insertnote\end{minipage}}
\fi
\tikzset{fase/.style={draw=none,fill=#1,rounded corners=3pt,minimum width=158pt,minimum height=118pt,text width=146pt,align=center}}

\renewcommand\ITERAmeeting{Pertemuan 5}
\begin{document}
\SlideJudul{IF25-11002 PRAKTIKUM PEMROGRAMAN}{Pertemuan 5\\Perulangan Bagian 1}{for, while, dan do-while dengan pola pencacah, akumulator, dan sentinel}{Tim Pengajar}{Tahun 2026. Sejajar dengan Pertemuan 5 Algoritma Pemrograman: Perulangan Bagian 1.}
\note{Buka dengan menanyakan satu hal dari pertemuan lalu yang masih membingungkan.}

\begin{frame}{Dua mata kuliah, satu minggu}
\framesubtitle{Sebelum mulai}
Topik minggu ini sama di dua mata kuliah, dengan peran yang berbeda.
\vspace{10pt}

\begin{columns}[T]
\begin{column}{0.47\textwidth}
\begin{Blok}{Algoritma: Perulangan Bagian 1}
Memodelkan pengulangan dengan flowchart dan pseudocode ULANGI dan SELAMA, lalu menelusurinya dengan trace table.
\end{Blok}
\end{column}
\begin{column}{0.47\textwidth}
\begin{BlokGelap}{Praktikum: Perulangan 1}
Menulis \texttt{for}, \texttt{while}, dan \texttt{do-while}, memakai pencacah dan akumulator, dan menghentikan perulangan dengan sentinel.
\end{BlokGelap}
\end{column}
\end{columns}
\note{Tanyakan apa yang sudah dibahas di Algoritma minggu ini. Gunakan jawabannya sebagai jembatan ke praktikum.}
\end{frame}

\begin{frame}{Saat pulang nanti, kalian bisa}
\framesubtitle{Target hari ini}
\begin{columns}[T]
\begin{column}{0.47\textwidth}
\begin{Langkah}{1}{Menerapkan}
Menulis program yang mengulang proses dengan \texttt{for}, \texttt{while}, dan \texttt{do-while}, memakai pola pencacah, akumulator, dan nilai sentinel

{\footnotesize\color{ITERAInkMuted}Sub-CPMK 5.1, CPMK0607}
\end{Langkah}
\end{column}
\begin{column}{0.47\textwidth}
\begin{Langkah}{2}{Menjelaskan}
Menelusuri perulangan dengan trace table, menentukan berapa kali badan perulangan dijalankan, dan menjelaskan penyebab off-by-one atau perulangan tak berhenti

{\footnotesize\color{ITERAInkMuted}Sub-CPMK 5.2, CPMK0608}
\end{Langkah}
\end{column}
\end{columns}
\vspace{8pt}
\begin{Catatan}
Dua kemampuan ini dinilai sama berat di Exit Ticket, tugas, UTS, dan UAS.
\end{Catatan}
\note{Bacakan target dengan pelan. Kata kuncinya menerapkan dan menjelaskan.}
\end{frame}

\begin{frame}{Ingat minggu lalu}
\framesubtitle{Review, 5 menit}
\begin{columns}[T]
\begin{column}{0.31\textwidth}
\begin{BlokGelap}{Pertanyaan 1}
Apa keluaran rantai \texttt{else if} untuk nilai batas 75?
\end{BlokGelap}
\end{column}
\begin{column}{0.31\textwidth}
\begin{BlokGelap}{Pertanyaan 2}
Mengapa perlu \texttt{break} di \texttt{switch}?
\end{BlokGelap}
\end{column}
\begin{column}{0.31\textwidth}
\begin{BlokGelap}{Pertanyaan 3}
Apa beda \texttt{if (x = 5)} dan \texttt{if (x == 5)}?
\end{BlokGelap}
\end{column}
\end{columns}
\vspace{8pt}
Jawab di kertas dulu, lalu cocokkan dengan teman sebelah.\note{Pilih dua mahasiswa untuk menjawab. Koreksi singkat, jangan mengajar ulang.}
\end{frame}

\Bagian{1}{Masalah pembuka}{Rata-rata nilai satu kelas}
\note{Masalah pembuka 10 menit.}

\begin{frame}[fragile]{Rata-rata nilai satu kelas}
\framesubtitle{Masalah minggu ini}
\begin{columns}[T]
\begin{column}{0.55\textwidth}
{Asisten praktikum harus menghitung rata-rata nilai Exit Ticket. Jumlah mahasiswa berbeda di setiap kelas, kadang 30, kadang 45.

\vspace{6pt}
Menulis \texttt{cin} sebanyak 45 kali jelas tidak masuk akal. Lagi pula banyaknya data baru diketahui saat program berjalan.\par}
\vspace{6pt}
\begin{Catatan}
\textbf{Diskusi 2 menit.} Langkah apa yang terus diulang untuk setiap nilai? Bagaimana program tahu kapan berhenti?
\end{Catatan}
\end{column}
\begin{column}{0.41\textwidth}
\begin{Keluaran}[]
Masukkan nilai (-1 untuk selesai):
70 85 60 90 -1
Banyak    : 4
Tertinggi : 90
Rata-rata : 76.25
\end{Keluaran}
\end{column}
\end{columns}
\note{Tampung beberapa jawaban tanpa menilai benar salah.}
\end{frame}

\begin{frame}{Dari langkah ke instruksi}
\framesubtitle{Sambungan dengan Algoritma}
\begin{columns}[T]
\begin{column}{0.31\textwidth}
\begin{Langkah}{1}{Temukan yang berulang}
Baca satu nilai, tambahkan ke total, tambah pencacah.
\end{Langkah}
\end{column}
\begin{column}{0.31\textwidth}
\begin{Langkah}{2}{Tentukan kapan berhenti}
Berhenti saat bertemu -1, atau setelah n kali.
\end{Langkah}
\end{column}
\begin{column}{0.31\textwidth}
\begin{Langkah}{3}{Terjemahkan}
\texttt{while} untuk syarat berhenti, \texttt{for} untuk banyak putaran yang sudah diketahui.
\end{Langkah}
\end{column}
\end{columns}
\vspace{10pt}
Langkah 1 dan 2 dilatih di Algoritma. Langkah 3 kita kerjakan hari ini dengan C++.\note{Hubungkan eksplisit ke Algoritma minggu ini.}
\end{frame}

\Bagian{2}{Coba dulu, baru dijelaskan}{25 menit eksperimen di GDB Online}
\note{Struggle 25 menit. Berkeliling, jangan menjelaskan materi dulu.}

\begin{frame}{Misi kalian}
\framesubtitle{Struggle, 25 menit}
\begin{columns}[T]
\begin{column}{0.31\textwidth}
\begin{Langkah}{1}{Hitung}
Tampilkan angka 1 sampai 10 tanpa menulis \texttt{cout} sepuluh kali.
\end{Langkah}
\end{column}
\begin{column}{0.31\textwidth}
\begin{Langkah}{2}{Jumlahkan}
Hitung 1 + 2 + ... + 100. Tebak dulu hasilnya.
\end{Langkah}
\end{column}
\begin{column}{0.31\textwidth}
\begin{Langkah}{3}{Berhenti}
Baca angka terus sampai pengguna mengetik 0.
\end{Langkah}
\end{column}
\end{columns}
\vspace{8pt}
\begin{columns}[T]
\begin{column}{0.47\textwidth}
\begin{Blok}{Boleh}
Bertanya ke teman, mengubah apa saja, sengaja membuat error.
\end{Blok}
\end{column}
\begin{column}{0.47\textwidth}
\begin{BlokAlert}{Belum dulu}
Mencari jawaban jadi di internet atau AI.
\end{BlokAlert}
\end{column}
\end{columns}
\note{Catat kesalahan yang sering muncul untuk dibahas di debrief.}
\end{frame}

\begin{frame}{Apa yang kalian temukan?}
\framesubtitle{Debrief struggle}
\begin{columns}[T]
\begin{column}{0.31\textwidth}
\begin{BlokGelap}{Pertanyaan 1}
Variabel apa yang kalian butuhkan untuk menyimpan jumlah?
\end{BlokGelap}
\end{column}
\begin{column}{0.31\textwidth}
\begin{BlokGelap}{Pertanyaan 2}
Di mana variabel itu harus diberi nilai awal?
\end{BlokGelap}
\end{column}
\begin{column}{0.31\textwidth}
\begin{BlokGelap}{Pertanyaan 3}
Apakah ada yang programnya tidak mau berhenti? Mengapa?
\end{BlokGelap}
\end{column}
\end{columns}
\vspace{10pt}
Jawaban kalian akan kita uji di bagian berikutnya.\note{Tulis jawaban di papan; buktikan atau patahkan di bagian materi.}
\end{frame}

\Bagian{3}{Mengulang dengan terkendali}{for, while, do-while, dan pola-pola perulangan}
\note{Materi 40 menit. Tunjukkan setiap contoh langsung di GDB Online.}

\begin{frame}[fragile]{Perulangan for}
\framesubtitle{Banyak putaran diketahui}
\begin{columns}[T]
\begin{column}{0.55\textwidth}
\begin{Kode}[]{for.cpp}
for (int i = 1; i <= 5; i++) {
    cout << i << " ";
}
cout << endl;
for (int i = 10; i > 0; i -= 3) {
    cout << i << " ";
}
cout << endl;
\end{Kode}
\end{column}
\begin{column}{0.41\textwidth}
\only<1>{\begin{TebakDulu}
{\ITERAKickerFont\fontsize{10pt}{12pt}\selectfont\color{ITERAAlert}TEBAK DULU}\par\vspace{4pt}
Sebelum dijalankan, tulis keluarannya di kertas.
\end{TebakDulu}
}\only<2>{\KeluaranFile[]{keluaran/P05/k001.txt}
\begin{Catatan}
Tiga bagian \texttt{for}: nilai awal; syarat lanjut; perubahan. Syarat diperiksa sebelum setiap putaran.
\end{Catatan}
}
\end{column}
\end{columns}
\note<1>{Minta mahasiswa menebak keluaran sebelum membuka slide berikutnya. Tanyakan satu atau dua tebakan yang berbeda, lalu buktikan.}
\end{frame}

\begin{frame}[fragile]{Pencacah dan akumulator}
\framesubtitle{Pola perulangan}
\begin{columns}[T]
\begin{column}{0.55\textwidth}
\begin{Kode}[]{akumulator.cpp}
int jumlah = 0, genap = 0;
for (int i = 1; i <= 10; i++) {
    jumlah += i;
    if (i % 2 == 0) {
        genap++;
    }
}
cout << "Jumlah: " << jumlah << endl;
cout << "Genap : " << genap << endl;
\end{Kode}
\end{column}
\begin{column}{0.41\textwidth}
\only<1>{\begin{TebakDulu}
{\ITERAKickerFont\fontsize{10pt}{12pt}\selectfont\color{ITERAAlert}TEBAK DULU}\par\vspace{4pt}
Sebelum dijalankan, tulis keluarannya di kertas.
\end{TebakDulu}
}\only<2>{\KeluaranFile[]{keluaran/P05/k002.txt}
\begin{Catatan}
Akumulator dan pencacah selalu diberi nilai awal sebelum perulangan, bukan di dalamnya.
\end{Catatan}
}
\end{column}
\end{columns}
\note<1>{Tanya: apa yang terjadi bila int jumlah = 0 dipindah ke dalam perulangan?}
\end{frame}

\begin{frame}[fragile]{Perulangan while}
\framesubtitle{Berhenti karena syarat}
\begin{columns}[T]
\begin{column}{0.60\textwidth}
\begin{Kode}[listing options={style=ITERACodeKecil}]{while.cpp}
int saldo = 100000, bulan = 0;
while (saldo < 500000) {
    saldo += 150000;
    bulan++;
}
cout << bulan << " bulan, saldo " << saldo << endl;
\end{Kode}
\end{column}
\begin{column}{0.36\textwidth}
\only<1>{\begin{TebakDulu}
{\ITERAKickerFont\fontsize{10pt}{12pt}\selectfont\color{ITERAAlert}TEBAK DULU}\par\vspace{4pt}
Sebelum dijalankan, tulis keluarannya di kertas.
\end{TebakDulu}
}\only<2>{\KeluaranFile[listing options={style=ITERAPlainKecil}]{keluaran/P05/k003.txt}
\begin{Catatan}
Sesuatu di dalam badan harus mengubah syarat. Bila tidak, perulangan tidak pernah berhenti.
\end{Catatan}
}
\end{column}
\end{columns}
\note<1>{Minta mahasiswa menebak keluaran sebelum membuka slide berikutnya. Tanyakan satu atau dua tebakan yang berbeda, lalu buktikan.}
\end{frame}

\begin{frame}[fragile]{Nilai sentinel}
\framesubtitle{Membaca sampai tanda berhenti}
\begin{columns}[T]
\begin{column}{0.60\textwidth}
\begin{Kode}[listing options={style=ITERACodeKecil}]{sentinel.cpp}
int nilai, banyak = 0, total = 0;
cin >> nilai;
while (nilai != -1) {
    banyak++;
    total += nilai;
    cin >> nilai;
}
cout << banyak << " data, total " << total << endl;
\end{Kode}
\end{column}
\begin{column}{0.36\textwidth}
\begin{Keluaran}[title={Masukan},listing options={style=ITERAPlainKecil}]
70 85 60 90 -1
\end{Keluaran}
\only<1>{\begin{TebakDulu}
{\ITERAKickerFont\fontsize{10pt}{12pt}\selectfont\color{ITERAAlert}TEBAK DULU}\par\vspace{4pt}
Sebelum dijalankan, tulis keluarannya di kertas.
\end{TebakDulu}
}\only<2>{\KeluaranFile[listing options={style=ITERAPlainKecil}]{keluaran/P05/k004.txt}
\begin{Catatan}
Baca sekali sebelum perulangan, lalu baca lagi di akhir badan. Dengan begitu -1 tidak ikut dihitung.
\end{Catatan}
}
\end{column}
\end{columns}
\note<1>{Minta mahasiswa menebak keluaran sebelum membuka slide berikutnya. Tanyakan satu atau dua tebakan yang berbeda, lalu buktikan.}
\end{frame}

\begin{frame}[fragile]{Perulangan do-while}
\framesubtitle{Minimal sekali jalan}
\begin{columns}[T]
\begin{column}{0.60\textwidth}
\begin{Kode}[listing options={style=ITERACodeKecil}]{dowhile.cpp}
int pin, coba = 0;
do {
    cin >> pin;
    coba++;
} while (pin != 1234 && coba < 3);
cout << (pin == 1234 ? "Masuk" : "Terkunci");
cout << " setelah " << coba << " kali" << endl;
\end{Kode}
\end{column}
\begin{column}{0.36\textwidth}
\begin{Keluaran}[title={Masukan},listing options={style=ITERAPlainKecil}]
1111 1234
\end{Keluaran}
\only<1>{\begin{TebakDulu}
{\ITERAKickerFont\fontsize{10pt}{12pt}\selectfont\color{ITERAAlert}TEBAK DULU}\par\vspace{4pt}
Sebelum dijalankan, tulis keluarannya di kertas.
\end{TebakDulu}
}\only<2>{\KeluaranFile[listing options={style=ITERAPlainKecil}]{keluaran/P05/k005.txt}
\begin{Catatan}
Syarat diperiksa sesudah badan. Jangan lupa titik koma setelah \texttt{while (...)}.
\end{Catatan}
}
\end{column}
\end{columns}
\note<1>{Minta mahasiswa menebak keluaran sebelum membuka slide berikutnya. Tanyakan satu atau dua tebakan yang berbeda, lalu buktikan.}
\end{frame}

\begin{frame}{Memilih jenis perulangan}
\framesubtitle{for, while, atau do-while}
\small
\begin{tabular}{@{}>{\raggedright\arraybackslash}p{364pt}>{\raggedright\arraybackslash}p{153pt}>{\raggedright\arraybackslash}p{345pt}@{}}
\kepala{Situasi} & \kepala{Pilihan} & \kepala{Contoh}\\
Banyak putaran diketahui & \texttt{for} & tampilkan 1 sampai n\\
\barisgenap Berhenti karena syarat, bisa nol putaran & \texttt{while} & baca sampai sentinel\\
Badan harus jalan minimal sekali & \texttt{do-while} & menu, validasi PIN\\
\garisdasar
\end{tabular}
\vspace{10pt}

Ketiganya bisa saling menggantikan. Pilih yang paling jelas menyatakan maksud kalian.
\end{frame}

\begin{frame}[fragile]{Tebak keluarannya}
\framesubtitle{Latihan menjelaskan, 3 menit}
\begin{columns}[T]
\begin{column}{0.56\textwidth}
\begin{Kode}[]{tebak.cpp}
int x = 0;
for (int i = 0; i < 4; i++) {
    x += i * 2;
}
int k = 1;
while (k < 20) {
    k *= 3;
}
cout << x << " " << k << endl;
\end{Kode}
\end{column}
\begin{column}{0.40\textwidth}
\begin{Catatan}
Tulis keluarannya di kertas, karakter demi karakter. Jangan dijalankan dulu.
\end{Catatan}
\end{column}
\end{columns}
\note{Contoh soal Bagian B. Beri waktu 3 menit sebelum slide jawaban.}
\end{frame}

\begin{frame}[fragile]{Tebak keluarannya: jawaban}
\framesubtitle{Latihan menjelaskan}
\begin{columns}[T]
\begin{column}{0.56\textwidth}
\begin{Kode}[]{tebak.cpp}
int x = 0;
for (int i = 0; i < 4; i++) {
    x += i * 2;
}
int k = 1;
while (k < 20) {
    k *= 3;
}
cout << x << " " << k << endl;
\end{Kode}
\end{column}
\begin{column}{0.40\textwidth}
\begin{Keluaran}[]
12 27
\end{Keluaran}
\begin{Catatan}
\texttt{x} = 0 + 2 + 4 + 6 = 12. \texttt{k} menjadi 3, 9, 27; berhenti saat 27 tidak lagi kurang dari 20.
\end{Catatan}
\end{column}
\end{columns}
\note{Tanyakan jawaban yang paling sering salah, lalu telusuri bersama.}
\end{frame}

\begin{frame}{Error yang sering muncul minggu ini}
\framesubtitle{Pesan diambil langsung dari g++}
\footnotesize
\begin{tabular}{@{}>{\raggedright\arraybackslash}p{208pt}>{\raggedright\arraybackslash}p{256pt}>{\raggedright\arraybackslash}p{398pt}@{}}
\kepala{Kesalahan} & \kepala{Contoh} & \kepala{Pesan pertama dari g++}\\
Koma di header for & \texttt{for (int i = 0, i < 5, i++)} & \texttt{error: expected ';' before '<' token}\\
\barisgenap Variabel di luar jangkauan & \texttt{cout << i; sesudah for} & \texttt{error: 'i' was not declared in this scope}\\
Lupa ; sesudah do-while & \texttt{\} while (x < 5)} & \texttt{error: expected ';' before 'cout'}\\
\barisgenap Perbandingan bertipe beda & \texttt{i < s.length() dengan int} & \texttt{warning: comparison of integer expressions of different signedness: 'int' and ...}\\
Badan for kosong & \texttt{for (...);} & \texttt{warning: this 'for' clause does not guard...}\\
\garisdasar
\end{tabular}
\vspace{10pt}

{\small Pesan di tabel ini dihasilkan dengan g++ dan opsi -Wall, sama seperti di cek.sh.}
\note{Minta mahasiswa memotret slide ini.}
\end{frame}

\Bagian{4}{Hands-on bertingkat}{45 menit, dari dasar sampai tantangan}
\note{Mahasiswa membuka folder 02 Hands-on/P5 - Perulangan 1 - Hands-on.}

\begin{frame}{Peta hands-on}
\framesubtitle{Folder 02 Hands-on/P5 - Perulangan 1 - Hands-on}
\small
\begin{tabular}{@{}>{\raggedright\arraybackslash}p{36pt}>{\raggedright\arraybackslash}p{267pt}>{\raggedright\arraybackslash}p{110pt}>{\raggedright\arraybackslash}p{83pt}>{\raggedright\arraybackslash}p{341pt}@{}}
\kepala{No} & \kepala{File} & \kepala{Tingkat} & \kepala{Waktu} & \kepala{Yang dilatih}\\
1 & \texttt{handson1-deret.cpp} & Dasar & 10' & \texttt{for}, akumulator\\
\barisgenap 2 & \texttt{handson2-tabungan.cpp} & Menengah & 15' & \texttt{while} dengan syarat berhenti\\
3 & \texttt{handson3-rerata.cpp} & Tantangan & 20' & trace table, perbaiki tiga kesalahan sentinel\\
\barisgenap + & \texttt{bonus-faktorial.cpp} & Pengayaan & sisa & \texttt{do-while} untuk menu, \texttt{long long}\\
\garisdasar
\end{tabular}
\vspace{10pt}

Kerjakan berurutan. Cocokkan keluaran dengan folder \texttt{expected/}; latihan yang membaca masukan memakai data di folder \texttt{input/}.\note{Saat berkeliling, minta mahasiswa yang mengerjakan latihan tantangan menjelaskan blok PENJELASAN mereka.}
\end{frame}

\begin{frame}{Cara mengecek dan aturannya}
\framesubtitle{Hands-on}
\begin{columns}[T]
\begin{column}{0.47\textwidth}
\begin{Blok}{Di GDB Online}
Salin isi file latihan, lengkapi TODO, klik Run. Bila program membaca masukan, ketik isi file di \texttt{input/}. Bandingkan dengan \texttt{expected/}.
\end{Blok}
\end{column}
\begin{column}{0.47\textwidth}
\begin{BlokGelap}{Di komputer sendiri}
Jalankan \texttt{bash cek.sh}. Skrip mengompilasi dengan \texttt{-Wall -Werror}, memberi masukan dari \texttt{input/}, lalu mencocokkan keluaran.
\end{BlokGelap}
\end{column}
\end{columns}
\vspace{6pt}
\begin{Catatan}
AI boleh untuk memahami pesan error, tidak untuk meminta solusi utuh. Exit Ticket dikerjakan tanpa bantuan apa pun.
\end{Catatan}
\note{Hands-on tidak dinilai langsung; yang dinilai Exit Ticket dan tugas.}
\end{frame}

\Bagian{5}{Exit Ticket dan tugas}{25 menit terakhir, lalu pekerjaan rumah}
\note{Mulai Exit Ticket tepat waktu.}

\begin{frame}{Exit Ticket 5}
\framesubtitle{Individual, 25 menit, tanpa AI dan internet}
\small
\begin{tabular}{@{}>{\raggedright\arraybackslash}p{60pt}>{\raggedright\arraybackslash}p{560pt}>{\raggedright\arraybackslash}p{80pt}>{\raggedright\arraybackslash}p{150pt}@{}}
\kepala{Soal} & \kepala{Isi} & \kepala{Poin} & \kepala{CPMK}\\
A1 & Tulis program yang menampilkan kelipatan k dari 1 sampai n beserta banyaknya & 25 & CPMK0607\\
\barisgenap A2 & Tulis program yang membaca bilangan sampai 0 lalu menampilkan jumlah bilangan positif & 25 & CPMK0607\\
B1 & Buat trace table sebuah perulangan dan tentukan keluarannya & 20 & CPMK0608\\
\barisgenap B2 & Jelaskan mengapa sebuah perulangan tidak berhenti dan cara memperbaikinya & 20 & CPMK0608\\
B3 & Tentukan berapa kali badan perulangan dijalankan dan jelaskan alasannya & 10 & CPMK0608\\
\garisdasar
\end{tabular}
\vspace{10pt}

{\small Soal A dinilai dari keluaran yang tepat sama. Soal B dinilai dari ketepatan dan alasan. Pelanggaran aturan membuat nilai Exit Ticket ini 0.}
\note{Soal lengkap dibagikan terpisah dan tidak disimpan di repo publik.}
\end{frame}

\begin{frame}{Tugas 5: Program Statistik Nilai}
\framesubtitle{Dikumpulkan sebelum Pertemuan 6}
\begin{columns}[T]
\begin{column}{0.47\textwidth}
\begin{Blok}{Bagian A, program, 50 poin}
Buat program statistik nilai dengan ketentuan berikut. Ketentuannya di slide berikut.
\end{Blok}
\end{column}
\begin{column}{0.47\textwidth}
\begin{BlokGelap}{Bagian B, penjelasan, 50 poin}
Komentar maksimal 150 kata dan trace table untuk n = 3 dengan nilai 80, 70, 90. Jelaskan alasan memilih \texttt{for}, \texttt{while}, atau \texttt{do-while} di setiap bagian, dan satu kesalahan off-by-one atau perulangan tak berhenti yang kalian temui.
\end{BlokGelap}
\end{column}
\end{columns}
\vspace{8pt}
{\small Data di tugas adalah data kalian sendiri. Rubrik lengkap ada di Modul Belajar 5.}
\note{Ingatkan bahwa penjelasan disalin dari AI mudah dikenali karena tidak menyebut pengalaman mereka sendiri.}
\end{frame}

\begin{frame}{Ketentuan Tugas 5}
\framesubtitle{Program Statistik Nilai}
\small
\begin{tabular}{@{}>{\raggedright\arraybackslash}p{87pt}>{\raggedright\arraybackslash}p{788pt}@{}}
\kepala{No} & \kepala{Ketentuan Bagian A}\\
1 & Baca jumlah siswa n, ulangi sampai n > 0\\
\barisgenap 2 & Baca nilai n siswa; setiap nilai harus 0 sampai 100, ulangi bila tidak valid\\
3 & Tampilkan nilai tertinggi, nilai terendah, dan rata-rata kelas\\
\barisgenap 4 & Tampilkan jumlah siswa LULUS (nilai 75 ke atas) dan TIDAK LULUS\\
5 & Menu \texttt{do-while} untuk mengulang atau keluar\\
\barisgenap Bonus & Tampilkan persentase kelulusan dengan dua desimal.\\
\garisdasar
\end{tabular}
\note{Bacakan ketentuan satu per satu dan tanyakan apakah ada yang belum jelas.}
\end{frame}

\begin{frame}{Rubrik Tugas 5}
\framesubtitle{Empat level, sama di semua instrumen}
\footnotesize
\begin{tabular}{@{}>{\raggedright\arraybackslash}p{166pt}>{\raggedright\arraybackslash}p{44pt}>{\raggedright\arraybackslash}p{166pt}>{\raggedright\arraybackslash}p{156pt}>{\raggedright\arraybackslash}p{146pt}>{\raggedright\arraybackslash}p{146pt}@{}}
\kepala{Kriteria} & \kepala{Poin} & \kepala{Sangat baik 85--100} & \kepala{Baik 70--84} & \kepala{Cukup 55--69} & \kepala{Kurang <55}\\
A1 Program berjalan & 20 & tanpa error dan peringatan & ada peringatan & perlu satu perbaikan & tidak terkompilasi\\
\barisgenap A2 Validasi dan perulangan & 15 & validasi n dan nilai, menu berulang & satu validasi kurang & dua kurang & tanpa validasi\\
A3 Ketepatan statistik & 15 & semua tepat termasuk n = 1 & satu salah & dua salah & sebagian besar salah\\
\barisgenap B1 Trace table dan alasan & 30 & trace lengkap dan alasan tepat & satu baris trace salah & trace tidak lengkap & tidak ada atau disalin\\
B2 Cerita error dan pengujian & 20 & pesan, sebab, perbaikan, uji & pesan tidak dikutip & hanya perbaikan & tidak ada\\
\garisdasar
\end{tabular}
\vspace{8pt}

{\small Nilai kriteria = poin × skor level. Bagian A total 50, Bagian B total 50.}
\note{Rubrik ini sama strukturnya setiap minggu; yang berubah hanya isi kriteria A2, A3, dan B1.}
\end{frame}

\begin{frame}{Sebelum Pertemuan 6}
\framesubtitle{Persiapan}
\begin{columns}[T]
\begin{column}{0.31\textwidth}
\begin{Langkah}{1}{Selesaikan}
Hands-on yang belum lulus \texttt{cek.sh}, terutama latihan tantangan.
\end{Langkah}
\end{column}
\begin{column}{0.31\textwidth}
\begin{Langkah}{2}{Kumpulkan}
Tugas 5 sebelum Pertemuan 6 dimulai.
\end{Langkah}
\end{column}
\begin{column}{0.31\textwidth}
\begin{Langkah}{3}{Baca}
Modul Belajar 5 untuk mengulang, lalu bagian awal modul berikutnya.
\end{Langkah}
\end{column}
\end{columns}
\vspace{10pt}
Pertemuan 6 membahas perulangan bersarang, \texttt{break}, \texttt{continue}, dan pola bintang, sejalan dengan Algoritma Pertemuan 6.\note{Tanyakan satu hal yang paling membingungkan hari ini untuk dibuka di review minggu depan.}
\end{frame}

\SlidePenutup{Terima kasih}{Sampai jumpa di Pertemuan 6: perulangan bersarang dan pola}
\note{Slide penutup.}
\end{document}
